% Pass 2 cuts (necessity audit) from Chapter 5 (Building adaptive hybrid representations).
% Each block below is the full notes block as it stood after pass 1; a shortened version replaced it in the chapter file.
% Cut: per-fold and per-model number lists (held-out evaluator pairs, stability rerun with one fold fewer, per-model transfer rows, decoy counts beyond one line);
% the section roadmaps that repeated subsection headings; literature detail beyond one citation per point (feature-selection pitfalls, PCA units and nonlinear methods, pairwiseMKL, Dablander p.138, Laufkoetter MCC, Turkina abstract);
% implementation detail not needed to follow a takeaway (rounding rule, trivial predictors for every metric, wide-representation pre-screen, DeepLIFT training settings, evaluator fallbacks, SVM cap counting on LogD/TATA, k-means cluster-edge test);
% repeats across sections (budget example in Results, evaluator evidence repeated in Results and Discussion, search-size question in Discussion and Conclusions, federated weighting formula);
% possible later additions (activity-cliff, developability and binding free-energy benchmarks); the word "donor", replaced by "augmentation set".

% ===== FROM: Introduction notes =====
% ---- NOTES: Introduction ----
% Takeaway: Standard representations carry many features that do not matter for a given task, and the best representation changes with the dataset and the model.
% Takeaway: This chapter builds the hybrid and reports when and why it helps. Distances and label noise are left to Chapter 6.
% Paragraph 1: the bridge from the earlier chapters
% - Chapter 4: standard representations carry many features that do not vary, or do not matter, on a given dataset.
% - A distance computed over all of those features inherits the problem.
% - Chapter 4 ended by asking for features chosen for the task from several representations, with irrelevant ones removed first.
% - Chapter 3: the best representation changed with the dataset and the model. That is the practical reason to search rather than fix one choice.
% - One clause only on the origin: the selection was first built for similarity, then found to work as model input.
% Paragraph 2: what this chapter does
% - It builds the adaptive hybrid representation ("hybrid") and reports first results on when and why it helps.
% - The idea gets one sentence here.
% Paragraph 3: scope
% - This chapter measures prediction only.
% - Distances in a hybrid, noise detection and learning with noisy labels are in Chapter 6.
% - The methods are final. Each dataset is split into 5 outer folds.
% - Result needed: on how many datasets the hybrid beats the best single representation, and whether the preferred search settings vary between datasets.
% Paragraph 4: roadmap, one sentence per part
% - Introduction: combining representations and the alternatives; the hybrid; blueprints; domains beyond small molecules.
% - Methods, the hybrid search: how a hybrid is built; the search space and why each setting is searched; augmentation; the fast tier and adaptive pruning.
% - Methods, experimental setup: datasets and splitting; representations by domain; models and metrics; comparators, controls and ablations; transfer, antibody scans and separate sites.
% - Results: when the hybrid helps; why it helps; blueprints; transfer and antibody scans; data held at separate sites. Then Discussion, Limitations and Conclusions.
% - Figure: chapter-opening overview. Left: several base representations as matrices of different widths (descriptors, fingerprint, graph embedding, sequence embedding). Middle: a few highlighted features from each. Right: the narrow hybrid going into any downstream model.
% ---- END NOTES ----

% ===== FROM: Combining representations and the alternatives notes =====
% ---- NOTES: Combining representations and the alternatives ----
% Takeaway: Most prior work joins whole representations or pools the predictions of one model per representation. The hybrid selects individual features across representations.
% Takeaway: No single representation wins across tasks, and new representations keep arriving, so the choice has to be made per task.
% Paragraph 1: prior work that combines fingerprints
% - Turkina et al. is the closest prior work. It combines and filters at once, but only across six non-hashed fingerprints.
% - Its abstract says bits are "selectively combined based on their importance within individual models". It does not say whether each fingerprint gets a fixed share.
% - TODO(lit): read the Turkina2025 full text for whether a fixed number of bits is kept per fingerprint, or all bits are ranked together.
% Paragraph 2: what this work adds
% - Change: "noise mitigation tasks" should read "noise detection and learning with noisy labels".
% - Change: "concatenate features" should read "select and join features". The hybrid joins features selected from each representation, not whole representations.
% - Augmentation gets one clause here.
% - "We explore using GNN embeddings" holds. MHG-GNN and GROVER are base representations on every small-molecule property dataset except QM9. GraphMVP is used on QM9 and logD. Name them, or cite Table~\ref{tab:reps_smallmol}.
% Paragraph 3: concatenation, late fusion, intermediate fusion, kernels
% - Change: the authors are Deng, Yang, Wang, Ojima, Samaras and Wang \citep{Deng2023}, not "Li and Deng".
% - In their study, appending RDKit2D descriptors to GROVER lowered RMSE on 9 of 24 datasets. Cite it for concatenation hiding which representation earned the score, not for concatenation degrading \citep{Deng2023}.
% - Multiple kernel learning: add its drug-bioactivity example, pairwiseMKL \citep{Cichonska2018}.
% - Closing sentence stays: no fusion strategy has been shown broadly superior \citep{Steyaert2023}.
% Paragraph 4: one representation by default, and the cost of joining several
% - Change: FernandezDelgado2014 compares classifiers on fixed datasets and never adds irrelevant features. It does not support "distance-based models are affected most severely".
% - Why distances suffer is argued in Chapter 4; one clause.
% - Your claim: a concatenation of two representations usually scores between its two parts, not above the better one. State it only if the two-representation concatenation test shows it.
% - Proposed reason: it adds the irrelevant features of both parts along with the useful ones. The method has to combine and filter at once.
% Paragraph 5: new representations keep arriving
% - What the method offers, in two sentences. A new embedding enters as one more base representation, with no change to the method. A new downstream model gets its own search, because models may want different features.
% - Change: typo "farmework".
% Paragraph 6: no single representation wins
% - This is the evidence that the best representation depends on the task. Chapter 3 showed the same on this thesis's own data.
% Paragraph 7: the range of representations (about two sentences)
% - Change: replace the domain-by-domain list with one sentence citing the representation tables (Tables~\ref{tab:reps_smallmol}, \ref{tab:reps_reactions}, \ref{tab:reps_antibodies}, \ref{tab:reps_nucleotides}, \ref{tab:reps_tcga}).
% - Keep: over 40 molecular foundation models were published between 2020 and 2024 \citep{Song2026}. Keep \citep{Bommasani2021} for models appearing faster than they are compared.
% - Change: the El Nahhas et al. sentence (ensemble of four pathology foundation models) is an ensembling result; it belongs in the ensembling paragraph.
% - Change: remove "presumably" and "somehow". Single modalities combined with independent ones gave more sensitive and robust predictions \citep{Steyaert2023}.
% Paragraph 8: where the hybrid sits among prior combination methods
% - Change: cut the first two sentences. They repeat Paragraph 3 (concatenation degrading; multiple kernel learning giving no interpretable selection).
% - Change: the consensus sentence \citep{Valsecchi2020} belongs in the ensembling paragraph.
% - Change: write "the hybrid search", not "KIRBy". Write "all base representations side by side", not "the full pool".
% - The hybrid works on individual features. It keeps original features drawn from several representations.
% - Closest framing in automated machine learning: choosing a learning algorithm and its hyperparameters together \citep{Thornton2013}. The hybrid search adds a third choice: which features, from which representations, go in.
% - Change: the information bottleneck sentences belong in the Discussion. Tishby2015 does not show that compression helps, so it cannot support them.
% Paragraph 9: ensembling
% - One model per representation; their predictions are pooled.
% - Consensus QSAR does this: it combines the predictions of several models, not their features \citep{Valsecchi2020}.
% - Second example: an ensemble of four pathology foundation models, 31 clinical tasks, mean AUROC 1.2\% above the best single model \citep{ElNahhas2025}.
% - No single model sees features from two representations together.
% - Every representation and every model must be run for each new molecule.
% Paragraph 10: feature selection, in three families
% - Feature selection keeps a subset of the original features. It builds no new ones. This separates it from PCA.
% - Filter methods score features against the labels before any model is fitted. Example: mRMR ranks features by mutual information with the label, minus redundancy with features already chosen \citep{Peng2005}.
% - Wrapper methods search over feature subsets, scoring each subset by fitting the model \citep{Kohavi1997}. In QSAR, descriptors were selected this way with a k-nearest-neighbour model \citep{Zheng2000}.
% - Embedded methods select while the model is fitted. Example: the lasso's L1 penalty sets some coefficients to exactly zero \citep{Tibshirani1996}.
% - Recursive feature elimination fits a model, drops the lowest-ranked features, and repeats \citep{Guyon2002}.
% - Reviews: \citep{Guyon2003}; for high-dimensional biological data \citep{Saeys2007}.
% - Two known pitfalls, one sentence each. Selected sets can change between resamples of the same data; the Nogueira index measures this across folds \citep{Nogueira2018}. Random-forest impurity importance favours some feature types over others \citep{Strobl2007}.
% - Where the hybrid sits: it ranks features by a model's importance scores (embedded), then keeps a fixed number (like a filter).
% - Classical methods treat all features as one flat list. They ignore which representation each feature came from.
% - Nearest classical relative: selection over predefined groups. The group lasso keeps or drops whole groups of variables together \citep{Yuan2006}.
% - The hybrid does not keep or drop whole representations. It gives each representation a share of the budget, then takes that many top-ranked features from it.
% - Because only original features are kept, every feature in the hybrid traces back to one representation. This is what makes blueprints possible.
% Paragraph 11: dimensionality reduction
% - It replaces the original features with fewer new, combined features.
% - In QSAR it turns many correlated descriptors into fewer uncorrelated ones \citep{Niazi2023}.
% - PCA is the standard method. Each component is a linear combination of the original features, ordered by the variance it explains \citep{Jolliffe2016}.
% - Components are usually functions of all original features, which makes them hard to interpret \citep{Jolliffe2016}.
% - PCA depends on the units of the features; z-scoring each feature first removes this \citep{Jolliffe2016}. This matters when fingerprint bits, descriptors and embeddings are joined.
% - PCA never sees the labels. It keeps the directions of largest variance, which need not predict the property. Components with small variance can still predict it \citep{Jolliffe1982}.
% - Partial least squares is the supervised counterpart in chemometrics: its components are built to covary with the property \citep{Wold2001}.
% - Both give mixed components. A component cannot be traced to one substructure, descriptor or representation, so no blueprint can be read from it.
% - Nonlinear methods, one sentence: kernel PCA \citep{Niazi2023}, t-SNE \citep{vanderMaaten2008}, UMAP \citep{McInnes2018}. t-SNE and UMAP are used mainly for two- or three-dimensional maps, not as model inputs.
% Paragraph 12: shrinking the feature space
% - Every alternative above except ensembling also reduces the number of features.
% - Models may want different numbers of features. Distance-based models (kNN, SVM) use every feature in every distance. Tree ensembles choose one feature at each split, so they can partly ignore irrelevant features.
% - Models may also want different features, not only different numbers.
% - Selection that uses the labels may overfit. Sort \& Slice, which uses no labels, beat two label-based substructure-selection schemes \citep{sns}. Dablander suggests the label use may add overfitting risk \citep[p.~138]{Dablander2024}.
% - Table: the alternatives side by side. One row per method: concatenation, late fusion and ensembling, intermediate fusion, multiple kernel learning, PCA and PLS, flat feature selection, group lasso, the hybrid. Columns: uses the labels; keeps original features; one model sees features from several representations; adding a representation needs retraining everything. Each row carries its citation.
% ---- END NOTES ----

% ===== FROM: The adaptive hybrid representation notes =====
% ---- NOTES: The adaptive hybrid representation ----
% Takeaway: A hybrid is up to a budget of features selected from several base representations of the same molecules, chosen by a search for one dataset and one downstream model.
% Takeaway: Filtering out useless features matters as much as picking useful ones, because pooled irrelevant features act as noise.
% Paragraph 1: no single representation is best for every dataset, even for one model
% - Change: name the method "adaptive hybrid representation", "hybrid" for short. Drop "KIRBy (Knowledge Integration by Representation)" and "framework".
% - Change: "the noise robustness study above" should read "the noise robustness study in Chapter~\ref{chap:robustness}".
% - Example model: NGBoost, Gaussian noise, from that study's own results.
% - QM9 clean R² on held-out molecules (median of 10 replicates): PDV 0.865 is highest, MHG-GNN 0.852 is close behind, ECFP4 0.706 is lowest.
% - logD clean R² (all 5 folds for every representation): PDV 0.673 is highest, MHG-GNN 0.658 second.
% - hERG Ki clean R²: MHG-GNN 0.511 is highest (all 5 folds). PDV 0.456 is second, the median of the 4 folds whose clean R² was at least 0.3.
% - So NGBoost's best representation changes with the dataset: PDV on QM9 and logD, MHG-GNN on hERG Ki.
% - Caco-2 (log efflux ratio): Avalon 0.402 and ChemBERTa 0.404 are medians of only 2 of 5 folds. The other 3 folds fell below clean R² 0.3 and were excluded.
% - Over all 5 folds, Caco-2 reverses. Avalon's median clean R² is 0.287, ChemBERTa's 0.256, and PDV's 0.370.
% - Change: the Caco-2 claim, that NGBoost prefers Avalon, holds only after 3 of Avalon's 5 folds are dropped.
% - Change: "only strong with PDV on QM9" overstates it. MHG-GNN is 0.013 R² lower.
% - Keep robustness apart from accuracy. On QM9, NGBoost's highest noise robustness (normalised area under the R²-noise curve) is on Avalon, 0.982 against 0.975 for PDV.
% - Change: cut the sentence guessing that Avalon captures features relevant to efflux. Nothing tested it.
% - Change: typos. "overidden" should be "overridden"; "its own type" should be "their own type"; "data set" is "dataset" elsewhere.
% Paragraph 2: what a hybrid is, in words
% - An adaptive hybrid representation is a feature set built by selecting up to a fixed number of features (the budget) from several base representations of the same molecules.
% - It is chosen by a search, for one dataset and one downstream model. A different dataset, split or model gives a different hybrid.
% - It is an input, not a model. Any standard supervised model can be trained on it.
% - Its intended use: representations that say different things about the molecule (PDV, ECFP4, graph embeddings, sequence embeddings), with the features for this task picked from all of them.
% Paragraph 3: why select (her second paragraph, plus one sentence)
% - Pooling every feature adds irrelevant features that act as noise (Chapter~\ref{chap:flaws}).
% - Add one sentence: distance-based noise handling did worst in an earlier exploratory study (Section~\ref{sec:lnl_motivation}).
% - That study is motivation, not a result of this thesis. It covered one dataset (QM9 HOMO-LUMO gap) and one noise type (Gaussian).
% - The label-noise question is measured in Chapter~\ref{chap:similarity}, over more datasets and noise types.
% - Change: typos. "the feature which carry" should be "the features which carry"; "data set" is "dataset".
% Paragraph 4: why the budget is split per representation, and what the search has to decide
% - Representations differ in width by orders of magnitude, so wide ones can crowd out narrow ones in a single ranking.
% - Importance scores from different feature types are not on a common scale.
% - The search also includes one pooled allocation method, so per-representation and pooled selection are compared, not assumed.
% - One sentence listing what the idea leaves open:
%   - how useful each feature is; how many features each representation gets; how many representations take part;
%   - how a candidate hybrid is judged during the search; the total budget; whether the answer depends on the downstream model.
% Paragraph 5: two ways to use a hybrid for similarity
% - One sentence: both are defined in Chapter~\ref{chap:similarity}, "Similarity in a hybrid".
% - Figure: several base representation matrices of different widths; within each, features ranked by importance; the budget as a bar split between representations; the top features from each joined into a narrow hybrid. Side panel: the full concatenation, for contrast.
% ---- END NOTES ----

% ===== FROM: Blueprints and what they tell a chemist notes =====
% ---- NOTES: Blueprints and what they tell a chemist ----
% Takeaway: A finished search is saved as a blueprint, and every selected feature traces back to its representation, so the search says something about the problem, not only its score.
% Takeaway: Named features let a chemist look for new chemistry or confirm the selection makes sense. Learned embedding features cannot be read this way.
% Paragraph 1: the output is more than a score
% - A finished search is saved as a blueprint: the settings needed to rebuild the hybrid.
% - Every selected feature traces back to its representation and its importance rank.
% - So the search returns information about the problem, not only a better score.
% - The search keeps the best configuration for each downstream model and budget, so one dataset gives one blueprint per model and budget. Comparing them at one budget shows whether tree ensembles and neural networks draw on different features.
% Paragraph 2: which features a chemist can read
% - What a chemist can read: descriptor names (e.g. molecular weight, logP), the atom environment behind a selected ECFP bit, the substructure behind a MACCS key.
% - Outside small molecules the readable unit changes: a sequence position and residue for antibodies, a pathway for gene expression.
% - Learned embedding features say which representation was useful, not which chemistry. Chapter~\ref{chap:flaws} covers why.
% Paragraph 3: two directions, discovery and confirmation
% - Discovery: substructures or descriptors that recur among the kept features point to parts of the molecule that may be involved in binding or activity. A starting point for a chemist, not a finding.
% - Confirmation: if the kept features make chemical sense for the endpoint, that supports the claim that the selection finds real signal and is not fitting noise.
% - Caveat: selected means predictive, not causal. Correlated features can stand in for each other, so an unselected feature may matter just as much.
% - Caveat: a blueprint rebuilt on other molecules can select other features. Read features from the fit on the dataset in question.
% - Context for explainability in drug discovery: \citep{JimenezLuna2020}. Interpretable inputs as an alternative to explaining a black box after the fact: \citep{Rudin2019}.
% Paragraph 4: other uses of a task-specific similarity
% - A hybrid also gives a distance chosen for one task. Anything that uses nearby molecules as evidence could use it instead of a standard fingerprint distance.
% - One sentence naming the uses argued in Chapter~\ref{chap:flaws}: noise detection, nearest-neighbour prediction, applicability domain, activity-cliff detection, clustering molecules for group-aware splits, and weighting predictions from separate sites.
% - One use not covered there: the predictor that scores molecules in generative design \citep{Olivecrona2017}. This is future work.
% ---- END NOTES ----

% ===== FROM: Beyond small molecules notes =====
% ---- NOTES: Beyond small molecules ----
% Takeaway: The search needs only a supervised task and several vector representations of the same samples, so it applies to tabular data in general.
% Takeaway: A dataset needs enough samples, and base representations whose features carry importance.
% Paragraph 1: what the method needs, and so where it applies
% - The search needs a supervised task and several vector representations of the same samples. Each representation is one feature matrix, one row per sample.
% - The default search objective scores prediction only. It reads no chemical structure.
% - So it applies to tabular data in general: any task where several encoders or measurements exist for the same inputs.
% - What does not carry over is chemistry-specific. Reading substructures from a blueprint needs named substructure features.
% - An optional search objective scores shared substructures between molecule pairs. It applies only to molecules.
% - The small-molecule distance and label-noise work of Chapter~\ref{chap:similarity} also does not carry over.
% Paragraph 2: one-hot encodings against learned sequence embeddings
% - One-hot sequence encodings are very sparse. At most positions the residue never changes across a dataset, so those features carry no importance.
% - One-hot can still work: it did on some antibody datasets.
% - Result needed: the antibody datasets where one-hot was among the selected representations, with the score in the dataset's metric.
% - Recommendation: prefer language-model sequence embeddings over raw one-hot encodings.
% Paragraph 3: the other domains used in this thesis, one or two sentences each
% - Antibodies: the most work after small molecules. Give this domain the most weight.
% - Chemical reactions (yield prediction): proof of concept. The representations are computed on reactants, products or their difference.
% - Nucleotide sequences (histone-mark and promoter tasks): proof of concept.
% - Cancer genomics (TCGA): each representation is a different measurement of the patient (expression, methylation, copy number, miRNA, protein, slide images), not a different encoder of one input. The search treats both cases the same way.
% Paragraph 4: what a dataset needs for the hybrid to help
% - Enough samples. The threshold (about 200 samples) waits on the full runs.
% - Base representations whose features carry importance. Where every representation already scores poorly, there are too few useful features to improve on.
% - Figure: one row per domain (small molecule, antibody, reaction, DNA sequence, patient sample). Left column: the raw input. Middle: several representation matrices. Right: the same search box, outputting a hybrid. Only the left column changes between rows.
% ---- END NOTES ----

% ===== FROM: Methods: the hybrid search notes =====
% ---- NOTES: Methods: the hybrid search ----
% Takeaway: The search builds and scores every combination of setting values inside cross-validation, keeps a winner for each downstream model and budget, and scores it once on an untouched outer fold.
% Paragraph 1: lead-in, 3-5 sentences, no definitions
% - The section describes how a hybrid is built from several base representations of the same samples.
% - The search is a grid: every combination of setting values is built and scored inside the cross-validation folds. Adaptive pruning, when switched on, drops some combinations after the first two folds.
% - A winner is kept for each pair of downstream model and budget. The budget is a maximum number of features; a hybrid can have fewer.
% - The reported comparison runs the whole search inside each outer training fold and scores the winner once on the untouched outer test fold.
% Paragraph 2: roadmap, one sentence per subsection
% - From base representations to a hybrid: the steps from base representations to one feature matrix, which molecules each step sees, and the blueprint that stores the settings.
% - The search space: each setting searched, its values, and why it is searched rather than fixed. Its parts cover feature importance scoring, feature allocation strategies and internal evaluation.
% - Augmentation: offering a finished hybrid, or one full representation, a few features from small or specialised feature sets. The extras are added on top of the budget or swapped in for the weakest features.
% - Fast tier and adaptive pruning: two ways to make the search cheaper. The fast tier uses shorter value lists, fixed in advance. Adaptive pruning drops values that lose a vote on the first two folds of this dataset.
% ---- END NOTES ----

% ===== FROM: From base representations to a hybrid notes =====
% ---- NOTES: From base representations to a hybrid ----
% Takeaway: Each representation's features are ranked by importance, the budget is split between representations by how well each predicts alone, and the top features from each are joined.
% Takeaway: Every step uses the training molecules of the fold only. The outer test molecules score only the finished hybrid.
% Paragraph 1: what the search takes in and gives back (her overview paragraph; fix it, do not add to it)
% - Inputs: $K$ base representations of the same $n$ molecules, $\mathbf{X}^{(k)} \in \mathbb{R}^{n\times d_k}$. Widths $d_k$ run from tens to thousands of features.
% - Also given: the labels $\mathbf{y}$, one group label per molecule (e.g. its scaffold), the dataset's metric, the task type, a downstream model and a list of budgets.
% - Output: a hybrid of at most $B$ features, where $B$ is the budget. The budget is a maximum, not a target.
% - Change: "at most $b$ columns" should read "at most $B$ features", matching the notation below.
% - Change: the clause on augmentation features outside the budget holds only for uncapped augmentation. Capped augmentation stays within the budget.
% - Change: "Blueprint" should be lower case.
% - Change: "according to baseline performance" is true for four of the five allocation methods. The global pool ranks the offered features together; say so in one clause.
% - Change: "sources" should read "representations", as in the rest of the chapter.
% Paragraph 2: preprocessing inside each fold (her paragraph, \citep{Ambroise2002, Vabalas2019})
% - Missing values: per-feature median of the training fold. Features missing on every training molecule are dropped.
% - Change: "Inputs from different sources" and "drop columns" should read "representations" and "features".
% - Every hybrid feature, binary bits included, is standardised to mean 0 and variance 1 with training-fold statistics.
% - Change: the exception is too broad. Binary features stay 0/1 only when the search scores a single representation to pick the comparator.
% - Single representations scored on the outer test fold are standardised like the hybrid.
% - A representation wider than 5,000 features is first cut, inside each fold, to its 500 features with the most LightGBM splits (200 if the dataset has under 1,000 molecules). This applies only to that internal single-representation scoring.
% - The ranking model is a 50-tree LightGBM of depth 5, fitted on the fold's training molecules. Features are ranked by how many times its trees split on them, not by gain.
% - So for such a representation the comparator is picked on its top 500 (or 200) features, then scored on the outer test fold at full width. The hybrid search always sees the full width.
% Paragraph 3: the folds, and which molecules each step sees
% - Outer folds: 5 by default, group-aware, so no group appears on both sides. Classification folds are also stratified by class.
% - A representation learned from the labels is re-trained inside each outer training fold, on training molecules only.
% - Inside each training fold, one group-aware internal split holds out about 20% of the molecules as whole groups. It is the first fold of a 5-fold grouped split, so it is deterministic.
% - On Koenig expression, whose training folds hold four groups, the internal split holds out one of the four.
% - Importance and each representation's standalone score come from this training fold only. The outer test molecules score only the finished hybrid.
% Paragraph 4: importance and standalone score per representation
% - Importance: the importance method gives each feature of representation $k$ one score $s^{(k)}_j \ge 0$. Scores rank features within one representation only.
% - Standalone score: the internal evaluator is fitted on representation $k$ alone and scored on the internal held-out part, giving $r_k$.
% - $r_k$ is rescaled to $\tilde r_k$ so that 0 means "no better than a trivial predictor". This way the budget split does not depend on which metric the dataset uses.
% - Rescaling per metric. Spearman, Pearson, MCC, $R^2$: $\max(0, r_k)$. AUROC and concordance index: $\max(0, 2(r_k - 0.5))$, chance at 0.5.
% - RMSE: $\max(0, 1 - \mathrm{RMSE}_k/\mathrm{RMSE}_0)$, where $\mathrm{RMSE}_0$ is predicting the training mean. MAE: the same, with the training median.
% - Accuracy and weighted F1: $\max(0, (r_k - r_0)/(1 - r_0))$, where $r_0$ is always predicting the majority training class.
% - The trivial predictor is scored on the same held-out molecules.
% - A representation whose fit fails gets weight 0 and is recorded as a failure.
% Paragraph 5: splitting the budget between representations
% - Optional representation count $N$: only the $N$ representations with the highest $\tilde r_k$ on this fold go on.
% - Each representation gets a weight $w_k = g(\tilde r_k)$. The allocation method sets $g$.
% - Exact share: $b^*_k = B\, w_k / \sum_l w_l$. Rounded by largest remainder: take $\min(\lfloor b^*_k \rfloor, d_k)$, then give out the leftover features one at a time, largest fractional part first, skipping full representations.
% - Result: $\sum_k b_k = \min(B, \sum_k d_k)$ and $b_k \le d_k$. A representation with $w_k > 0$ keeps at least one feature. If every weight is 0, the budget is split equally.
% - The global pool does not use this rounding: each representation offers candidates, and one joint ranking keeps at most $B$.
% Paragraph 6: selecting features and assembling the hybrid
% - $S_k$ = the top $b_k$ features of representation $k$ by $s^{(k)}_j$, among those with $s^{(k)}_j > 0$. The same cut at 0 applies to the features offered to the global pool and to its final ranking.
% - If no feature of a representation scores above 0, its full $b_k$ is taken.
% - Why the cut: tree importance gives exactly 0 to every feature no tree split on. Filling the budget from those would pick features by their position in the matrix.
% - Hybrid: $\mathbf{H} = [\mathbf{X}^{(1)}_{:,S_1}, \dots, \mathbf{X}^{(K)}_{:,S_K}]$, with $\sum_k |S_k| \le B$. So a hybrid can be narrower than its budget.
% - Then $\mathbf{H}$ is preprocessed inside the fold, as in Paragraph 2.
% Paragraph 7: scoring configurations, and the reported comparison
% - A configuration is one value on every searched setting.
% - Each configuration: fit the downstream model on the hybrid's training molecules, score on the held-out molecules, average over folds.
% - The winner is picked separately for each pair of downstream model and budget, never across models.
% - That score chose among many configurations, so it is optimistic \citep{Varma2006, Cawley2010}.
% - The reported comparison is nested: the whole search runs inside each outer training fold, the winning settings are rebuilt on that fold, and the result is scored once on the untouched outer test fold.
% - Result needed: how much higher the inner search score is than the outer test score, in each dataset's own metric.
% Paragraph 8: two variants used in Chapter 6 (one or two sentences)
% - Under label noise, every training-side quantity uses the noisy labels. Held-out scoring uses clean labels unless stated.
% - The nearest-neighbour objective replaces both the standalone score and the downstream-model score with a nearest-neighbour score: each molecule's activity is predicted from its nearest neighbours in the space. No downstream model or internal evaluator is used.
% - Figure: the search pipeline. Left: several matrices of different widths. Middle, inside one outer training fold: the internal 80/20 group split; importance bars per representation; standalone scores feeding the budget split; top features pulled from each matrix. Right: the hybrid, the downstream model, the held-out score. Around it: the loop over configurations and "best per downstream model and budget". Mark where the blueprint is saved. Mark the outer test fold as never seen by the search. Inset: outer folds, each containing the whole search.
% ---- END NOTES ----

% ===== FROM: Blueprints notes =====
% ---- NOTES: Blueprints ----
% Takeaway: A blueprint stores the recipe, not the features. Rebuilt on other labelled molecules it can pick other features; for unlabelled molecules, the features chosen by a labelled fit are reused.
% Paragraph 1: what a blueprint is and what it stores (her first paragraph; fix it)
% - A blueprint is the recipe: the search settings needed to rebuild the winning hybrid. Every scored configuration has one. The search returns the winner's blueprint with its result.
% - Change: her list lacks the objective, the metric's direction and the augmentation settings (how many extra features, from which candidates, and the rule for admitting one). Add them in one clause.
% - Change: her list includes the step size. It is stored, but no searched allocation method reads it. Drop it from the list.
% - Stored for inspection only: each representation's best score across the downstream models, how many features each representation got in each fold, and which augmentation features were kept in each fold. These show which representations fed the hybrid and whether that holds across folds.
% - That stored score is not the standalone score that splits the budget. It is measured on every molecule the search saw, so a rebuild never reads it.
% - Change: the recipe does not recreate "this exact hybrid representation on any data". On other molecules a rebuild re-scores the representations and re-picks features.
% Paragraph 2: two ways to reuse a blueprint (her second paragraph; fix it)
% - Change: the blueprint stores no selected base-representation features, as her prose says. It does store the per-fold feature counts and the kept augmentation features.
% - Re-run the recipe on new labelled data: importance and standalone scores are recomputed, so the features adapt to the new dataset. It costs a new selection and can give a different feature set.
% - Reuse the exact features from a labelled fit: take the same features from the new molecules. This needs no labels, so it is the only option for unlabelled molecules. It assumes the same features matter on the new data.
% - Re-running the recipe without labels is refused. Every importance score would be 0 and the pick would fall back to feature order.
% - Result needed: whether re-running the recipe or reusing the exact features transfers better from a large dataset to a smaller related one.
% - Importance and allocation can also be computed separately in k-means clusters of the fingerprints, to ask whether the budget split changes across chemical space. One sentence.
% Paragraph 3: naming the selected features (her third paragraph; fix it)
% - Change: the names come from a separate explanation step run on a finished hybrid. The blueprint itself does not store them. Change "This blueprint will tell you" accordingly.
% - ECFP (Morgan) bits: named by the atom environment that sets the bit, written as a SMILES fragment, not SMARTS. It is taken from the first molecule, of at most 50 supplied, that has the bit set.
% - MACCS keys: named by their SMARTS pattern and a short name. The descriptor vector's features carry the descriptor names.
% - Outside small molecules: protein composition, descriptor and k-mer features; antibody composition, one-hot and BLOSUM positions, structural and inverse-folding scores; gene-expression pathway sets; PCA components.
% - Change: her prose names only learned embeddings as unreadable. Add the others. Bit index only: MAP4 and DRFP. Representation name and dimension index only: Avalon, Mordred, atom-pair, graph-kernel features and every learned embedding.
% - Figure: two paths. Left, "re-run the recipe": new labelled data, score representations, split the budget, pick top features, hybrid. Right, "reuse the exact features": stored features taken from unlabelled molecules, hybrid. Mark which path needs labels.
% ---- END NOTES ----

% ===== FROM: The search space and why each axis is swept notes =====
% ---- NOTES: The search space and why each axis is swept ----
% Takeaway: Six settings are searched as a full grid: downstream model, budget, representation count, importance method, allocation method and internal evaluator.
% Takeaway: Each is searched because different datasets and models prefer different values, and the settings interact.
% Paragraph 1: motivation
% - Change: her paragraph "When I began to implement this idea..." fits here, as the opening of this subsection. It now sits at the end of "Search configuration and internal evaluation".
% - Her three questions map onto settings of the search. How to measure importance is the importance method.
% - How many features to take from each representation is the allocation method, the budget and the representation count. How to score a feature set is the internal evaluator.
% - The answer to each question became a setting that is searched, per dataset.
% Paragraph 2: why a grid
% - Every combination of values is built and scored. Nothing is fixed in advance except the values each setting may take.
% - The claim of this chapter is that no single setting suits every dataset. A grid measures that instead of assuming it.
% - One line of motivation from the literature: no learning method is best on every problem \citep{Wolpert1996}. Motivation only, not proof.
% - A full product, not one setting at a time, because the settings interact: the best budget depends on the downstream model (example under Paragraph 5).
% Paragraph 3: the settings, in one table
% - Six settings are searched: downstream model, budget, representation count, importance method, allocation method, internal evaluator.
% - Their values in the full search and the fast tier are in Table~\ref{tab:search_tiers}. The defaults and the full count breakdown are in Additional files, Section~\ref{app:searchspace}.
% - Importance method: LightGBM gain, TreeSHAP or DeepLIFT.
% - Why the importance method is searched: importance is a property of a model, not of the data alone \citep{Dong2020}. The three methods come from different model families and rank features differently, and a different ranking builds a different hybrid.
% - Allocation method: how the budget is split among the chosen representations. Five methods.
% - Internal evaluator: matched (the downstream model itself), random forest, LightGBM or DNN.
% - Not searched inside one search: the objective (prediction by default; nearest-neighbour prediction in Chapter 6) and augmentation (none, uncapped, capped). Each is a separate search.
% Paragraph 4: the downstream model
% - The model trained on the finished hybrid. Five models, plus a Gaussian process on small datasets.
% - Why it is searched: the model decides what counts as a useful feature and how many features it can use. So the best hybrid for one model need not be the best for another.
% Paragraph 5: the budget
% - The most features the hybrid may have. Eight values, from 10 to 2,000 features.
% - Why it is searched: too few features leave informative ones out; too many add uninformative ones and room to overfit on small datasets. Budgets above the number of training molecules are kept.
% - Example (preliminary; to be replaced by the full runs), nanobody thermostability, 5 folds: SVM did best at 300 features (0.6384 Spearman). DNN, LightGBM and random forest did best at 50 features, XGBoost at 100.
% Paragraph 6: the representation count
% - Each dataset has a fixed set of base representations, chosen before the search (Table~\ref{tab:dataset_reps}). Fourteen of the 16 datasets have 10; H3K4me3 has 8 and TATA promoter 9.
% - The count sets how many of those may contribute features to the hybrid: 2, 3, 5 or all.
% - It is the only step that drops whole representations before allocation. Correlation between representations plays no part.
% - Why it is searched: pooled feature selection over all representations has no such setting, and datasets differ in how many representations are worth using. Every allocation method is paired with every count.
% - Some datasets do best with 2 or 3 representations, others with all 10.
% ---- END NOTES ----

% ===== FROM: Feature importance scoring notes =====
% ---- NOTES: Feature importance scoring ----
% Takeaway: Three importance methods from different model families rank features differently, and which one builds the best hybrid depends on the dataset and the downstream model.
% Paragraph 1: what an importance score is and what the search does with it
% - Within one representation, features differ in how much they say about the label. The search ranks them before choosing any.
% - An importance method gives every feature of one representation one non-negative score.
% - Each representation gets its own freshly fitted model, on its own features only. The downstream model is not used.
% - Only the order within a representation is used. Scores are never compared across representations or across methods, so no normalisation is applied.
% - Scores are computed inside each fold, on the internal training part only (80% of that fold's training molecules, split by group).
% - Before fitting, missing values are set to 0 and each feature is standardised to mean 0 and variance 1 on the same molecules.
% - Scores are computed once per fold and method, then reused by every allocation method, budget, internal evaluator and downstream model.
% - Change: "Tree-based feature importance measures favor features with high split gain" does not match \citep{Strobl2007}. That paper shows split-based importance is biased toward features with many possible split points. This matters because binary fingerprint bits sit beside continuous embeddings.
% - \citep{Ancona2018} supports the DeepLIFT sentence.
% Paragraph 2: the three methods, and how to name them in text
% - Names in text, tables and figures: "LightGBM gain", "TreeSHAP", "DeepLIFT" (Deep Learning Important FeaTures) \citep{Shrikumar2017}. Change: the prose spells it "DeepLift".
% - Define each once at first use: "LightGBM gain importance", "SHAP values for tree ensembles (TreeSHAP)", "DeepLIFT".
% - Family labels: split gain (LightGBM gain), Shapley values (TreeSHAP), reference-based attribution (DeepLIFT).
% - Change: "explanation-based" does not separate TreeSHAP from the others. TreeSHAP also reads trees, and DeepLIFT is also an explanation method. Use the family labels.
% - Call DeepLIFT's model "a small neural network with two hidden layers", to keep it apart from the DNN downstream model.
% - LightGBM gain: how much a feature reduced training loss where trees split on it \citep{ke2017lightgbm}. Model: 50 gradient-boosted trees, maximum depth 6. Score: total loss reduction of every split on the feature, summed over all trees; 0 if no tree splits on it.
% - TreeSHAP: how much a feature moves each molecule's prediction, as an exact Shapley value for tree models \citep{Lundberg2020}. Model: random forest, 50 trees, maximum depth 10 \citep{Breiman2001}. Score: mean absolute Shapley value over the training molecules.
% - DeepLIFT: how much a feature moves a small network's output away from its output on a reference input \citep{Shrikumar2017, Kokhlikyan2020}. Network: two ReLU hidden layers of 64 and 32 units, 100 full-batch epochs of Adam, learning rate 0.001 \citep{Kingma2015}. Score: mean absolute attribution over the training molecules, Rescale rule.
% - DeepLIFT's reference is all zeros, which after standardising is every feature at its training mean. For a fingerprint bit, that is the fraction of training molecules that have the bit.
% - Classification labels: TreeSHAP and DeepLIFT use the positive class for binary tasks and average absolute values over classes for multiclass tasks.
% - Pre-screen on wide representations: for LightGBM gain and TreeSHAP, only the 2,000 highest-variance features (variance measured before standardising) are fitted; the rest score 0. DeepLIFT scores every feature.
% - Share of features scoring above 0 on Caco-2, median over 10 representations and 3 folds: 30% for LightGBM gain, 70% for TreeSHAP, 90% for DeepLIFT (preliminary).
% - Table: one row per importance method. Columns: model fitted per representation; score of one feature; classification labels; representations over 2,000 features.
% Paragraph 3: the ranking comparison
% - Two comparisons were run. The first compared how methods rank features; the second compared the hybrids those rankings build.
% - The ranking comparison scored 8 importance methods on 7 datasets (LogD, hERG Ki, Buchwald-Hartwig yield, nanobody thermostability, antibody FcRn retention, BRCA PAM50 subtype, LUAD vs LUSC). It measured the Spearman correlation between two methods' scores over every feature of every representation (preliminary).
% - The three kept methods rank features differently: Spearman 0.20-0.64 (LightGBM gain vs TreeSHAP), 0.21-0.67 (LightGBM gain vs DeepLIFT) and 0.56-0.99 (TreeSHAP vs DeepLIFT) across the 7 datasets.
% - Caveat: each method was fitted once on the whole dataset, not inside folds.
% - Change: the prose says the methods' "predictions did not correlate". It is their feature rankings that were compared; say "rankings".
% - Change: "across data sets and model architectures" needs both comparisons and the win counts behind it: rankings on 7 datasets, hybrids on one dataset with two downstream models, win counts on 5 datasets with 5 downstream models.
% Paragraph 4: the hybrid comparison
% - Hybrids were built on Caco-2 (2,158 molecules, 10 representations; 3 scaffold folds; budget 100 features; random forest and LightGBM downstream) from each method's ranking. Three control rankings were included: feature variance, Laplacian score and random.
% - Measure: RMSE of the best single representation minus RMSE of the hybrid, in log10 efflux ratio units. Positive means the hybrid has the lower error.
% - Result (preliminary): DeepLIFT +0.0144, LightGBM gain +0.0100, TreeSHAP +0.0067, random ranking -0.0118. The gap between DeepLIFT and LightGBM gain (0.0044 RMSE) is smaller than the fold-to-fold spread.
% - "Did not correlate strongly" fits the top of the rankings on Caco-2: of each representation's 100 top-ranked features, the kept pairs share a median of 28-54 features.
% Paragraph 5: the selection rule
% - Rule 1: a method must score every feature of every representation, for regression, binary and multiclass labels. All three kept methods do.
% - Rule 2: of two methods whose rankings agree above Spearman 0.95, keep the faster one.
% - Under rule 2 DeepLIFT replaced integrated gradients (above 0.95 on 4 of 6 datasets) and gradient SHAP (above 0.95 on 3 of 6).
% - Rule 3: keep methods that rank features differently, because a different ranking builds a different hybrid.
% - TreeSHAP came last of the three on Caco-2 and was kept under rule 3.
% - The only filters on features before allocation are the cut at a score of 0 and the pre-screen on representations over 2,000 features.
% - What the pre-screen drops on a 2,048-bit fingerprint: the 48 lowest-variance bits. On hERG Ki all 48 are among the 105 constant ECFP4 bits. On QM9, which has no constant ECFP4 bits, they are near-constant ones.
% - Near-constant features otherwise stay candidates and are ranked by importance.
% Paragraph 6: no single method suits every dataset and model
% - Record (preliminary): 5 datasets (nanobody thermostability, hERG Ki, humanness, LogD, BRCA PAM50 subtype); 5 folds; budgets 25, 50, 100, 150 and 300 features; 10 representations; 5 downstream models.
% - For each (downstream model, budget) pair, 25 per dataset, count which importance method gave the best configuration.
% - The pattern: which method gives the best configuration depends on the dataset and the downstream model. LightGBM gain won most pairs on four datasets; DeepLIFT won most on one.
% - LogD is that one: DeepLIFT won 17 of its 25 pairs. Outside LogD, DeepLIFT gave the best configuration at 2 of 100 pairs.
% - Change: "Networks preferred tree-based measures" holds on four of five datasets. With the DNN downstream, LightGBM gain won 17 of 20 budgets outside LogD. On LogD the DNN took DeepLIFT at 3 of 5 budgets.
% - Change: "Tree ensembles preferred gradient-based approaches" holds on LogD only. There XGBoost and LightGBM took DeepLIFT at 5 of 5 budgets each; random forest took LightGBM gain at 5 of 5.
% - Table: one row per dataset. A win is one (downstream model, budget) pair where that method gave the best configuration, out of 25. Best score is the best configuration reached with that method, in the dataset's own metric (preliminary).
%   Columns: dataset | metric | wins: LightGBM gain, TreeSHAP, DeepLIFT | best score: LightGBM gain, TreeSHAP, DeepLIFT
%   Nanobody thermostability | Spearman | 15, 10, 0 | 0.6696, 0.6471, 0.6440
%   hERG Ki | RMSE (pKi), lower is better | 14, 11, 0 | 0.5589, 0.5605, 0.5728
%   Humanness | AUROC | 12, 12, 1 | 0.9723, 0.9713, 0.9689
%   LogD | RMSE (log units), lower is better | 8, 0, 17 | 0.4250, 0.4345, 0.4170
%   BRCA PAM50 subtype | weighted F1 | 21, 3, 1 | 0.9097, 0.9018, 0.9038
% - Result needed: the same win counts from the full runs.
% Paragraph 7: how the importance method enters the search
% - The full search sweeps all three methods. The fast tier uses LightGBM gain only.
% - When adaptive pruning is on (it is off by default), importance methods that rarely give the best configurations on the first two folds are dropped from the rest; at least one is always kept.
% ---- END NOTES ----

% ===== FROM: Feature allocation strategies notes =====
% ---- NOTES: Feature allocation strategies ----
% Takeaway: An allocation method decides how many of the budget's features each representation contributes. Four give fixed shares from each representation's standalone score; the fifth lets candidates compete in one pooled ranking.
% Takeaway: No allocation method was best in most comparisons, so all five are searched.
% Paragraph 1: what an allocation method decides, and the five that are searched
% - An allocation method decides how many of the budget's $B$ features each representation contributes. The importance method then decides which ones.
% - Five methods are searched: linear, squared, softmax, rank decay, and a global pool.
% - The first four give each representation a fixed share of the budget, computed from its standalone score.
% - The global pool lets candidate features offered by every representation compete in one importance ranking.
% - Change: under the global pool, not all features compete. Only the candidates each representation offers do.
% - Change: "select the top features from each source and concatenate" is true for the four fixed-share methods only.
% Paragraph 2: the steps every method shares
% - Change: the standalone score comes from the internal validation part of the training fold, not "the training set".
% - Allocation runs only over the representations the representation count keeps.
% - Cost: one fit per representation per fold and internal evaluator. Every budget, allocation method and importance method reuses these scores.
% - Change: the numbered list in the prose (fit the evaluator, rescale, split, take the top features, join) repeats "From base representations to a hybrid". Replace it with one sentence.
% - Change: the list's formula uses $n_r$ where the rest of the chapter uses $b_k$.
% Paragraph 3: the four fixed-share methods
% - They differ only in the weight $w_k$ computed from the rescaled score $\tilde r_k$. Their settings are fixed for every dataset, not searched.
% - Linear: $w_k = \tilde r_k$. A representation gets 0 features when its rescaled score is 0.
% - Squared: $w_k = \tilde r_k^{\,2}$ (power 2). A representation gets 0 features when its rescaled score is 0.
% - Softmax: $w_k = \exp\big((\tilde r_k - \max_l \tilde r_l)/T\big)$, $T = 1$. No representation gets 0 features because of its score.
% - Rank decay: $w_k = e^{-0.3\rho_k}$, where $\rho_k$ is the rank by score (0 for the best). No representation gets 0 features because of its score.
% - With 10 representations, under softmax the best takes at most 23.2% of the budget, because every rescaled score lies between 0 and 1.
% - Under rank decay only the order counts, so the best gets 27.3% and the tenth 1.8%, whatever the gaps.
% - Worked example: budget 100 features, three representations with rescaled scores 0.6, 0.5 and 0.2. Linear 46/39/15; squared 55/39/6; softmax 39/35/26; rank decay 44/32/24 features.
% - Change: "variants of performance-weighted allocation ... proportional to performance" fits linear only. Say "fixed-share methods" for the four.
% - Table: one row per method. Columns: weight $w_k$; fixed setting; how features are chosen; when a representation gets 0 features.
% - Table: the worked example, one row per fixed-share method, one column per representation.
% Paragraph 4: the global pool
% - Squared weights $\tilde r_k^{\,2}$ size each representation's offer, as shares of a pool of $2B$ candidate features.
% - An offer is capped at the representation's width in features, and only features scoring above 0 are offered.
% - The offered features are joined into one matrix. The same importance method is fitted once on that matrix, on the internal training part.
% - The hybrid keeps the top $B$ features of that joint ranking, again only those scoring above 0.
% - Each representation's final count is an outcome of the joint ranking, not set in advance. There is no one-feature floor. A representation gets 0 features when its score is 0 or none of its candidates reach the joint top $B$.
% - A feature's rank here depends on which features the other representations offered. Under fixed shares it does not.
% - Cost: one extra importance fit per configuration, on the pooled candidates.
% Paragraph 5: when a representation gets no features
% - Under all five, a representation also gets none when the representation count leaves it out or its standalone fit fails. If every weight is 0, the budget is split equally and the run reports this.
% Paragraph 6: what fixed shares assume
% - A strong representation may owe its score to a few features, so a large share can fill with weak ones.
% - Fixed shares cannot see complementarity. Two representations carrying the same information both get a share.
% - Of the five, only the global pool ranks features from different representations against each other. State this as a difference in mechanism, without a verdict.
% - Change: her two sentences on this say the same thing twice ("It is worth noting that this is not always the case..." and "it relies on the assumption that many features in a well-performing representation are valuable"). Merge them, and keep the complementarity clause.
% Paragraph 7: greedy forward selection, tried first and cut (four or five sentences)
% - Why it was tried: it picks features for how well they work together, which fixed shares cannot \citep{Kohavi1997, Guyon2003}.
% - What it did: each round it added the next 10 features from whichever representation most improved the internal validation score. It stopped when the budget was spent or after 3 rounds with no gain.
% - What it cost: up to $(B/10) \times K$ fits per configuration and fold for $K$ representations. That is 500 fits at 500 features and 10 representations, against 10 for the score-based methods.
% - What it scored (preliminary): on Jain antibody FcRn retention, 0.3600 Spearman against 0.4526 for linear, 16th of 17 methods. On breast cancer PAM50 subtype, 0.8551 accuracy against 0.8938 for linear, also 16th of 17.
% - It was removed from the search.
% - Change: drop the commented sentence that it "works better on larger datasets". The PAM50 set has 821 tumours and greedy still placed 16th of 17.
% Paragraph 8: how the five were chosen (run before the main experiments; four sentences)
% - Compared: 8 candidate allocation methods, each at 4 representation counts, on 9 regression datasets with 3 downstream models and 2 budgets. That is 54 (dataset, model, budget) combinations (preliminary).
% - A fixed-share method gave the single best strategy in 38 of the 54 combinations, the global pool in 16.
% - All five stay because no single method was best in most combinations.
% - Caveat: this comparison ran before the current score rescaling existed. It sized weights with an inner 3-fold group split, not the 80/20 internal split used now.
% ---- END NOTES ----

% ===== FROM: Search configuration and internal evaluation notes =====
% ---- NOTES: Search configuration and internal evaluation ----
% Takeaway: The internal evaluator scores each representation on its own during the search. That score splits the budget and decides which representations take part.
% Takeaway: The evaluator that best serves a downstream model is not always the same model family, so four evaluators are searched, matching included.
% Paragraph 1: what an internal evaluator is and where it acts
% - A model fitted inside the search on one representation at a time. It is separate from the downstream model scored at the end.
% - Features are standardised on the internal training part. The evaluator is fitted there and scored on the internal validation part, in the dataset's own metric.
% - That score sets each representation's share of the budget. It also ranks representations for the representation count.
% - It does not rank features inside a representation. The importance method does that.
% - Judging a feature set by a model's score is the wrapper idea \citep{Kohavi1997}.
% - Change: "a model whose inductive bias best fits the downstream model" is out of date. No evaluator is chosen to fit; all four are searched and the best configuration wins.
% - Change: kNN is not a downstream model in the search. Use SVM or DNN in the opening example.
% - Change: "as we found in our feature importance scoring analysis" names no result. Cite a result, or drop the clause.
% - Change: "the final downstream evalutors" should be "the downstream model"; "another axes" should be "another setting".
% Paragraph 2: the four searched values and what matched means
% - Four values are searched: matched, random forest, LightGBM and DNN. The fast tier searches the same four.
% - Matched means an evaluator of the same model family as the downstream model.
% - For a random forest, LightGBM or DNN downstream model, matched is already in the list, so those three are searched with 3 evaluators.
% - For an XGBoost or SVM downstream model, matched adds an XGBoost or SVM evaluator, so those two are searched with 4.
% - Where a Gaussian process is a downstream model, matched adds a Gaussian process evaluator, so it is also searched with 4.
% - Same family is not the same settings. Only the DNN evaluator shares the downstream DNN's architecture and training settings.
% - One difference remains for the DNN: its early-stopping 20% is drawn at random. The downstream DNN splits its 20% by the same groups as the outer split.
% - The tree evaluators are smaller: 50 trees with a depth limit, against 100 trees downstream. This keeps scoring every representation cheap.
% - The SVM and Gaussian process evaluators are fitted on at most 500 training molecules, because their cost grows with the square to the cube of the molecule count.
% - The random forest evaluator tries all features at each split on regression datasets, and sqrt(number of features) on classification. The downstream forest tries sqrt(number of features) on both.
% - If the DNN evaluator fails to fit on a fold, a random forest scores that fold and the run records a warning.
% - Every evaluator setting, beside the same model's settings as a downstream model, is in Table~\ref{tab:app:model_settings} (Additional files).
% Paragraph 3: a downstream model the search does not know (one or two sentences)
% - A user can pass their own downstream model, under any name, as long as it can be fitted and can predict.
% - Matched knows only a fixed list of model families. An unknown model gets no matched evaluator, and no warning is given.
% - That model is scored only by the other evaluators: by default random forest, LightGBM and DNN. If matched was the only evaluator requested, a random forest scores the representations.
% - A hand-written blueprint that says matched, rebuilt with a model the search does not know, is scored by LightGBM instead.
% - State it as a limit: a user with their own model should name an evaluator.
% Paragraph 4: why the evaluator is searched, with the evidence
% - Variable importance depends on the model \citep{Dong2020}. The representation a forest scores highest need not be the one an SVM uses best.
% - Matching is the obvious default. But the evaluator only has to rank representations, and another model may rank them better for the downstream model.
% - Matched is one of the four searched values, so the search already covers it. A mismatch winning is informative, not a negative result.
% - In the search (preliminary): 5 datasets (antibody thermostability (Tresanco), hERG Ki, humanness, LogD, breast cancer PAM50 subtype), 10 representations, budgets of 25 to 300 features, 5 folds. Scores are the search's own fold scores, not held out.
% - The matched evaluator gave the best score for its downstream model in 9 of 25 (dataset, downstream model) pairs.
% - Per budget, it was best in 43 of 125 (dataset, downstream model, budget) combinations. With 3 or 4 evaluators and no effect, about 38 of 125 would be expected.
% - Antibody thermostability, DNN downstream: matched 0.6696 Spearman against 0.6506 for the LightGBM evaluator, a gap of 0.0190. Matched's 5 fold scores ran from 0.5095 to 0.7920.
% - LogD, SVM downstream: matched 0.4322 RMSE (log units) against 0.4217 for the LightGBM evaluator (lower is better), a gap of 0.0105. Matched lost at all 5 budgets. Its 5 fold scores ran from 0.4118 to 0.4634.
% - Held out (preliminary): 3 outer folds, budget 250 features, 5 representations, LightGBM gain only. Random forest and SVM downstream.
% - Antibody thermostability, random forest: matched 0.6098 Spearman, LightGBM evaluator 0.6220, SVM evaluator 0.6172. Matched's fold scores ran from 0.5659 to 0.6734.
% - Antibody thermostability, SVM: matched 0.6189 Spearman, Gaussian process evaluator 0.6265, LightGBM evaluator 0.6317. Matched's fold scores ran from 0.5530 to 0.6703.
% - hERG Ki, random forest: matched 0.6046 RMSE (pKi), SVM evaluator 0.6059, LightGBM evaluator 0.6075. Matched's fold scores ran from 0.5957 to 0.6178. Matched beat the SVM evaluator on 2 of 3 folds.
% - hERG Ki, SVM: matched 0.5723 RMSE (pKi), LightGBM evaluator 0.5724, Gaussian process evaluator 0.5765. Matched's fold scores ran from 0.5605 to 0.5862.
% - Every gap above, held out and in the search, is smaller than the range of matched's own fold scores. State the pairs and the ranges; no verdict beyond the numbers.
% - Result needed: per dataset and downstream model, how often the winning configuration used the matched evaluator, and its gap to the best other evaluator in the dataset's metric.
% Paragraph 5: how the setting values were chosen (her paragraphs "When I began..." and "We ran a series of experiments...")
% - Before the full runs, smaller design runs fixed the values each setting may take: 1 to 9 datasets, 1 to 5 folds (preliminary). Each dataset was scored in its own metric; scores were never averaged across datasets.
% - What was measured: whether a value gave the best score for some dataset, downstream model and budget, and how many hours it cost.
% - A value that won somewhere but cost many hours stayed in the full search and left the fast tier.
% - Each setting's values were either kept, because each won for some dataset, model and budget, or fixed at one value, because the spread between values was smaller than one value's spread across folds.
% - Kept, example (preliminary): nine downstream models were fitted on each of 10 base representations, 5 folds, on Buchwald--Hartwig, Caco-2, Koenig expression, LogD and BRCA PAM50 subtype. The DNN scored best on 4 of the 5 and third on Caco-2 (0.2936 against 0.2851 RMSE, log$_{10}$ efflux ratio), so it was added.
% - Fixed, example (preliminary): rescaling each representation's score before the budget is split. On hERG Ki, linear allocation, 300 features, the five rescalings gave 0.5870 to 0.5924 RMSE (pKi). One of them alone ran 0.5475 to 0.6294 RMSE across the 5 folds. So rescaling is one fixed rule per metric.
% - The full list of design runs is the table in Additional files, Section~\ref{app:searchspace}.
% Paragraph 6: size of the search (her paragraph "The full search tests...")
% - 8,160 configurations per dataset holds where all five downstream models run. It is lower where a dataset drops a model.
% - The breakdown by downstream model and the number of model fits is in Additional files, Section~\ref{app:searchspace}.
% Paragraph 7: caching (her paragraph "The complete search result is cached...")
% - Importance scores and each representation's own score are computed once per fold and reused by every configuration. So cost grows more slowly than the configuration count.
% - A fast search can be extended to a full one.
% - The chosen configuration per (downstream model, budget) is saved as a blueprint.
% ---- END NOTES ----

% ===== FROM: Augmentation notes =====
% ---- NOTES: Augmentation ----
% Takeaway: Small or specialised feature sets cannot compete fairly in the budget split, so they are offered afterwards as extras, ranked by how well each explains what the base gets wrong.
% Takeaway: Augmentation also lets a chemist keep one trusted representation and learn which few outside features help, and how many.
% Paragraph 1: why some feature sets are handled apart from the budget split (her first paragraph)
% - Some feature sets are too small to compete for budget: MACCS keys are 167 features, ESM-2 log-likelihood ratios are 3 features.
% - Others only mean something next to a full representation, such as mutation features (differences from a parent antibody).
% - Leaving them out loses signal they may carry. Putting them in the budget split distorts it, because a 3-feature set cannot fill a large share.
% - So they are offered after the hybrid $\mathbf{H}$ is built. Call $\mathbf{H}$ the base, and the feature sets the extras are drawn from the donors.
% - Change: the prose describes only the uncapped mode, where extras are appended. The capped mode, where extras replace base features, is missing.
% Paragraph 2: which feature sets are donors (her second paragraph)
% - Donors are set per dataset, not per domain: the Augmentation column of Table~\ref{tab:dataset_reps}.
% - MACCS is a donor on four of the six small-molecule datasets. On MDR1-MDCK it is a base representation. Word it as "usually a donor".
% - QM9 is the exception: MACCS is not offered, and the topological torsion fingerprint (2048 features) is a donor. It is the one wide donor.
% - Two datasets have no donors, humanness and trastuzumab--HER2. The nanobody thermostability set has one, mutation features.
% - Change: the Coulomb matrix donor is not a fingerprint. It is the 50 largest eigenvalues of a matrix built from atomic numbers (nuclear charges) and 3D interatomic distances.
% - Change: spelling. "Coloumb" should be "Coulomb"; "sames" should be "samples" or "variants".
% - Common thread: each donor is small (at most 167 features, torsions aside) or describes one aspect only: shape, charge, the mutated positions, one tumour property. A plausible supplement and a weak stand-alone representation.
% - Why 3D donors are worth offering: \citet{Bahia2023} modelled six protein--ligand series (461 structures) with 2D, 3D and 2D+3D descriptors, using k-nearest neighbours, random forest and lasso. Combining 2D and 3D descriptors gave many more significant models. The authors attribute this to the two coding for complementary properties.
% - Difference to state if citing it: their 3D descriptors came from bioactive conformations in the PDB. Here the 3D donors use one generated conformer per molecule.
% Paragraph 3: ranking donor features by what the base gets wrong
% - Inside each outer training fold, an 80/20 internal split is made. Scaffold or sequence groups stay on one side. The outer test fold is never used.
% - A small random forest (50 trees, depth 10) is fitted on the base features, with 3-fold group-aware cross-validation on the internal-train molecules.
% - Its out-of-fold residuals are $r_i = y_i - \hat{y}^{\mathrm{oof}}_i$. For binary tasks, $\hat{y}$ is the predicted probability of the positive class.
% - Out-of-fold, because a forest's in-sample residuals shrink towards zero and mostly measure fit noise.
% - All donor features are pooled and each is scored as $c_j = |\rho_{\mathrm{P}}(\mathbf{x}_j, \mathbf{r})|$, the absolute Pearson correlation. The top $M = 50$ features go on.
% - Why residuals and not labels: ranking against $\mathbf{y}$ favours features that repeat what the base already knows. Ranking against $\mathbf{r}$ favours features that explain its errors.
% - Same idea as fitting each new stage to the current residual in gradient boosting \citep{friedman2001greedy}, and as least angle regression adding the variable most correlated with the residual \citep{efron2004lars}.
% - The residual model is always this cheap forest, never the downstream model. The ranking then costs the same for every downstream model.
% Paragraph 4: choosing how many extras to keep
% - Block sizes $k \in \{0, 5, 10, 20, 40\}$ are tried, capped at the number of candidates. Block $k$ is the top $k$ ranked features.
% - Each block is scored with the downstream model on the internal-validation molecules. The chosen size $k^*$ has the best score, and must beat $k=0$.
% - $k=0$ (no extras) is scored first and wins ties. "No extras" is the default outcome, not a failure.
% - Five whole-block choices replace about 50 one-by-one accept or reject decisions.
% - Choosing among many options on a small validation set overfits the choice \citep{Cawley2010}. Five options leave less room for that.
% - The chosen donor and feature positions are stored with the rest of the search settings. They are re-extracted for the outer fold like the rest of the hybrid.
% Paragraph 5: two ways to fit the extras against the budget $B$
% - Uncapped: extras are appended outside the budget. The final width is $|\mathbf{H}| + k^*$, at most $B + 40$ features.
% - Capped: extras must displace base features. The base features are ranked once by random forest importance, and the $k^*$ weakest are dropped. The final width stays at $B$.
% - In capped mode, any unspent budget is filled first, and at least one base feature is always kept.
% - Only in capped mode does an extra feature have to beat a base feature to earn its place.
% - Why both: uncapped and no augmentation at the same nominal $B$ differ in width. That comparison mixes "better features" with "more features". Capped removes the width difference.
% Paragraph 6: where augmentation enters the search, and the second use
% - Augmentation is a separate condition, not a searched setting. For each dataset with donors, the whole search runs three times: no augmentation, uncapped, capped.
% - Every other setting (downstream model, budget, allocation method, importance method, internal evaluator) is searched inside each condition.
% - 14 of the 16 datasets run all three conditions. Humanness and trastuzumab--HER2 run no augmentation only.
% - No augmentation: width at most $B$, nothing offered. Uncapped: width at most $B + 40$, an extra must beat keeping no extras. Capped: width at most $B$, an extra must beat keeping no extras and the weakest base feature it replaces.
% - Under the nearest-neighbour search objective (Chapter 6), all three conditions are also run.
% - Two datasets never run the nearest-neighbour objective: BRCA PAM50 subtype (4 classes; that objective is scored by AUROC, which needs 2) and trastuzumab--HER2 binding (a held-out block can hold one group only).
% - Second use, outside the main search: keep one strong representation at full width as the base and add a few extras. There $B$ is that representation's width, so only uncapped makes sense.
% - This use is tested in the experiment that adds a few features to one strong representation, with decoy features planted to check that the rule takes none.
% - Why it matters for adoption: a chemist can keep the representation they already trust and learn which few outside features help, and how many.
% - Augmentation asks directly whether another representation holds signal the base lacks (Chapter 4).
% - On Koenig expression, the mutation features and log-likelihood ratios reach the model only through augmentation.
% Paragraph 7: limitations
% - $|\rho_{\mathrm{P}}|$ sees only linear, one-feature-at-a-time relations with the residual.
% - It does not see redundancy among the extras. Two near-identical donor features can both rank high.
% - Donor features are pooled with no per-donor limit. A 2048-feature donor gets more chances in the ranking than a 49-feature one.
% - The margin over $k=0$ is zero, so a near-tie on one internal split can admit extras.
% - Table: the three conditions. One row per condition. Columns: final width (features); what an extra feature must beat; datasets that run it.
% - Figure: flow diagram. Left: base matrix $\mathbf{H}$ and donor blocks. Middle: forest on $\mathbf{H}$, out-of-fold residuals $\mathbf{r}$, donor features ranked by $|\rho_{\mathrm{P}}(\mathbf{x}_j,\mathbf{r})|$, top 50. Right: $k^*$ chosen from $\{0,5,10,20,40\}$ on internal validation, then two branches: uncapped appends $k^*$ features; capped drops the $k^*$ weakest base features and adds $k^*$.
% - Result needed: per dataset and downstream model, how often $k^*=0$ is chosen, and per-fold scores for no augmentation, uncapped and capped.
% - Result needed: which donors are kept most often, per domain.
% - Result needed: how many decoy features are admitted in the decoy check.
% ---- END NOTES ----

% ===== FROM: Fast tier and adaptive pruning notes =====
% ---- NOTES: Fast tier and adaptive pruning ----
% Takeaway: The fast tier searches shorter value lists, fixed in advance, and accepts a possibly worse answer. Adaptive pruning drops values that clearly lose on the first two folds and aims to find the same winner.
% Paragraph 1: two ways to make the search cheaper, and how they differ
% - The fast tier uses shorter lists on some settings. The lists are fixed before the search and are the same on every fast-tier dataset.
% - Adaptive pruning decides during the search, from this dataset's own fold scores.
% - The fast tier knowingly accepts a possibly worse answer in exchange for a much smaller grid. Pruning aims to find the same best configuration while no longer scoring configurations that have already clearly lost.
% - Without either, the full search does not finish on the largest datasets.
% Paragraph 2: which datasets run the fast tier, and what that costs
% - Six datasets run the fast tier (marked in Table~\ref{tab:dataset_reps}); all others run the full search.
% - Why: the full search would take longer than four 12-hour jobs in a row (about 40 hours), or the dataset has more than 10,000 molecules or sequences. Estimated hours at full: QM9 about 103, Koenig expression about 64, LogD about 62. H3K4me3 has 18,244 sequences.
% - The dropped values are still searched on every other dataset. The fast tier is a compute compromise, not a claim that those values are useless.
% - A fast run can be extended to a full one later without redoing work. Each configuration's score is stored by its settings, not by the tier, so a full search scores only the configurations the fast one skipped.
% Paragraph 3: how the fast lists were chosen (preliminary)
% - Rule: a value stays in the fast lists if it is cheap and adds to the best score. An expensive value that wins often stays in the full search only.
% - Importance methods and downstream models were measured on five datasets (TCGA breast cancer PAM50 subtype, nanobody thermostability, antibody humanness, hERG Ki, LogD), 5 folds each, against the unpruned full search.
% - DeepLIFT took about half of the fast run's hours on all five datasets (49.9% to 51.6%). Dropping it changed only LogD's best score, from 0.4905 to 0.4930 RMSE.
% - Adding DNN as a downstream model moved LogD from 0.4930 to 0.4273 RMSE (full-search best 0.4170). Adding SVM moved hERG Ki from 0.5776 to 0.5599 RMSE. Adding XGBoost changed no dataset's best score.
% - The global pool won 60 of the 125 combinations of dataset, downstream model and budget on those five datasets. It is expensive, so it is in the full search only.
% - The three fast allocation methods are linear, squared and rank decay.
% - Result needed: the share of run time the fast tier and pruning saved on the full runs.
% ---- END NOTES ----

% ===== FROM: Adaptive pruning notes =====
% ---- NOTES: Adaptive pruning ----
% Takeaway: Pruning scores every configuration on two folds, drops values that are rarely among the best within groups that fix every other setting, and never removes every value of a setting.
% Takeaway: Its cut-offs were set on five datasets so that no best score changed.
% Paragraph 1: what pruning is for
% - Some configurations are so far behind after two folds that the other three folds would only confirm a loss. Pruning stops scoring those.
% - It is judged on one criterion: it must never remove a configuration that would have been the best. Time saved is secondary.
% - Pruning is off unless a search asks for it. As submitted, the search stored once per dataset uses pruning on every dataset, fast and full tier alike.
% - On a fast-tier dataset only one importance method is searched, so the importance vote does not run.
% - Related idea in hyperparameter search: drop configurations early on partial evaluations, as in successive halving \citep{Jamieson2016} and Hyperband \citep{Li2018Hyperband}. Here the partial evaluation is two of the five folds.
% Paragraph 2: the steps
% - Change: her paragraph says "the first two folds"; the rules rank configurations by their mean score over folds 1 and 2, in the dataset's own metric.
% - The checkpoint at a budget of 100 features runs once every budget up to 100 has been scored. Budgets are scored smallest first.
% - The checkpoint ranks configurations at budgets up to 100 within each group of (downstream model, internal evaluator, representation count). A pair of allocation and importance method with no configuration in the better half is not scored above 100 features. At least one pair per group is kept.
% - Each vote groups configurations by every other setting, including the budget. The best few in each group get a point. A value is removed if it gets a point in fewer than the cut-off share of groups.
% - Representation count: best 1, cut-off 10% of groups. Allocation method: best 2, 20%. Importance method: best 2, 20%. Internal evaluator: best 3, 10%.
% - A vote runs only if the setting has more values than get a point. It can never remove every value of a setting.
% - A configuration that fails to fit is ranked last in its group.
% - Survivors are scored on all five folds. The best configuration is picked by the mean over five folds, exactly as without pruning.
% - Table: one row per rule. Columns: how groups are formed; who is kept or gets a point; cut-off share of groups; what it removed on the calibration datasets.
% Paragraph 3: why the votes are grouped, and why two folds
% - Grouping by every other setting means values compete only with each other. A value that is best for a random forest at 25 features is not removed for losing to a DNN at 500 features.
% - One fold is not enough. In the representation-count vote on TCGA breast cancer PAM50 subtype after fold 1 alone, 726 of 1,020 groups had every value at the identical score. Ties are then broken by list order, so a one-fold vote removes values by accident.
% - After one fold, no setting of any vote was safe on the calibration datasets. After two folds, each setting had at least one safe cut-off.
% Paragraph 4: how the cut-offs were set (preliminary)
% - Calibrated on five datasets: logD, hERG Ki, TCGA breast cancer PAM50 subtype, antibody thermostability (Tresanco) and antibody humanness.
% - Each dataset's grid was first scored unpruned on all five folds. Each rule was then replayed on the stored fold scores, so no rule needed new fits.
% - A cut-off counts as safe if the best five-fold score is unchanged for each of 25 combinations of downstream model and budget, on every dataset.
% - Where several safe cut-offs made the same cuts, the one with a safe neighbour was chosen. The saving was not the criterion.
% - What the chosen rules removed: 5 representations on logD and 2 on PAM50 subtype; softmax allocation on logD; DeepLIFT on antibody thermostability and hERG Ki, TreeSHAP on logD; the SVM evaluator on logD. The checkpoint removed configurations within groups only, never a whole allocation or importance method.
% - TreeSHAP is 71% of logD's fitting time and was never the best on logD. Most of the saving comes from there.
% - Estimated run time for eight budgets, unpruned against pruned: logD 79 h against 44 h, hERG Ki 42 h against 35 h, PAM50 subtype 42 h against 30 h.
% - The calibration applied the votes one after another. The search applies all four at once, on the same fold-1-and-2 scores. So the cut-offs were checked in a different order from the one the search uses.
% - Five datasets is a small calibration set. State this as a limit.
% Paragraph 5: what pruning does to the reported score
% - With pruning on, folds 1 and 2 both choose the survivors and score them. The five-fold mean of the winner is then slightly optimistic \citep{Cawley2010}.
% - Single representations are never pruned. A fair comparison of hybrid against single representation leaves both voting folds out of both, or re-scores without them.
% - Figure: a funnel. Top: every configuration on folds 1 and 2. Next: the checkpoint at 100 features. Next: the four votes. Bottom: survivors on folds 3 to 5, then the best configuration by the mean over all five folds. Label each stage with the share of configurations left on one calibration dataset.
% ---- END NOTES ----

% ===== FROM: Methods: experimental setup notes =====
% ---- NOTES: Methods: experimental setup ----
% Takeaway: Every experiment reuses one protocol: 5 outer group folds, the search on training molecules only, and a paired comparison with the best single representation on the same held-out molecules.
% Paragraph 1: what this section sets up, in the order of its subsections
% - Datasets and splitting: 16 datasets across 6 domains, and how their folds are cut.
% - Representations by domain: the base representations offered to the search for each domain.
% - Models and metrics: the downstream models, and the metric per dataset.
% - Comparators, controls and ablations: the experiments that test whether the hybrid's gain is real.
% - Transfer across datasets, antibody mutational scans and data held at separate sites: the setups for those three studies.
% Paragraph 2: the shared protocol
% - Each of the 16 datasets is split into 5 outer folds. No group is split across folds.
% - One exception: the antibody expression scan used 3 outer folds of unequal size. State the 3 where it is reported.
% - The hybrid search sees only the training molecules of the outer fold.
% - The selected features are rebuilt on those training molecules and applied unchanged to the held-out molecules.
% Paragraph 3: what the hybrid is compared against
% - The comparator is the best single base representation, used at its full number of features.
% - It asks whether the hybrid beats what a user would otherwise run.
% - It is chosen from the search's own single-representation scores on the training molecules, then refit and scored on the same held-out molecules.
% - It can be a different representation in each fold.
% - Choosing it on the held-out molecules would give it an advantage the hybrid does not get.
% - The hybrid is held to its budget of features. The comparator is not, so it is often many times wider. Both widths are recorded per fold; say so, so a difference is not read as purely a method effect.
% - The main quantity per fold is hybrid score minus best-single score, signed so that positive always means the hybrid scored better. For RMSE the sign is flipped.
% - Reported per dataset and model: the per-fold scores of the hybrid and of every single representation, and their spread.
% Paragraph 4: link to Chapter 6 (one or two sentences)
% - Chapter 6 reuses this protocol with noise added to the training labels of each outer fold. Held-out labels stay clean.
% ---- END NOTES ----

% ===== FROM: Models and metrics notes =====
% ---- NOTES: Models and metrics ----
% Takeaway: Five downstream models from different families score every hybrid and every single representation, with fixed settings, so differences come from the input.
% Takeaway: Each dataset is scored in its field's own metric and read on its own; results are per-fold scores and their spread, with no significance tests.
% Paragraph 1: the downstream models
% - Five downstream models score every hybrid and every single representation: random forest, XGBoost, LightGBM, SVM and DNN.
% - Together they cover bagged trees, boosted trees, a kernel method and a neural network.
% - A Gaussian process is added as a sixth model on datasets under 1,000 molecules \citep{Rasmussen2006}. Its kernel is RBF plus a white-noise term.
% - Three of the 16 datasets are under 1,000 molecules: humanness (564), nanobody thermostability (672) and breast cancer PAM50 (821). The Gaussian process runs on humanness only.
% - PAM50: left out under every augmentation setting for run time. It took 44 of the 60 hours of one search, the one with augmentation on top of the budget.
% - Nanobody thermostability: left out under augmentation on top of the budget, where its searches did not finish in the time limit. A dataset that drops the Gaussian process under any setting gets it under none, so it is left out everywhere.
% - Change: "Eleven models are available" and the list with CatBoost, NGBoost, the MC-Dropout BNN and the 1D CNN. None of these four is used in any experiment. Name only the models that are used.
% - Change: kNN is not a downstream model. Drop it from the opening sentence's examples.
% - Change: "model pool" should read "set of downstream models".
% - Change: "We use the full set for model-matching and noise experiments" no longer holds. An experiment that uses a model outside the five names it in its own methods.
% - Why five: the best set of features depends on the model that uses it. A hybrid that helped only one model family would be a property of that model, not a better input.
% - Tree ensembles are the baseline to beat: a hybrid is a tabular feature matrix, and tree-based models still lead on medium-sized tabular data of about 10,000 samples \citep{Grinsztajn2022}.
% Paragraph 2: fixed settings and size caps
% - Every model uses the same settings on every dataset, so differences come from the input.
% - There is no per-dataset tuning by default. An optional per-fold random search exists (20 draws, 3 group-aware inner folds, outer training fold only). The DNN is never tuned this way; it relies on early stopping.
% - Result needed: whether tuning pays on base representations (tuned against untuned models). One sentence, in each dataset's metric.
% - Gaussian process: skipped above 1,000 molecules. It builds and factorises the full molecule-by-molecule kernel.
% - SVM: skipped above 10,000 molecules. Above that its kernel cache no longer holds the support vectors and a fit stops finishing in useful time.
% - SVM under augmentation on top of the budget: skipped above 5,000 molecules.
% - The caps count the molecules the search is given. The search run once per dataset is given the whole dataset. A search inside one outer fold is given that fold's training molecules, about 80% of the dataset.
% - On two datasets the two counts give different downstream models, both under augmentation on top of the budget. LogD has 5,028 molecules and TATA promoter 5,274. The SVM is dropped when the whole dataset is counted, and kept when about 4,000 training molecules are counted.
% - Some default models are not run on named datasets where a single fit was too slow. They are listed in the Additional files table of downstream models not run.
% - Every downstream model's settings, library, molecule cap and citation are in Table~\ref{tab:app:model_settings} (Additional files).
% Paragraph 3: k-nearest neighbours, where it enters
% - kNN is not a default downstream model. Its default is $k = 5$ neighbours.
% - It is the model in the comparison of importance-weighted against unweighted distances, which reads 5, 10 and 20 neighbours.
% - State its settings once, here, because Chapter 7 "Simple is often better" relies on that result.
% Paragraph 4: metrics
% - Each dataset is scored in the metric its field normally uses, so its numbers can be compared with published work on the same dataset.
% - RMSE for small-molecule, quantum-chemistry and reaction-yield regression. It is in the label's own units.
% - Spearman $\rho$ for antibody regression. These assays are used to rank candidates, and their scales differ between labs.
% - AUROC for binary classification. It does not depend on a decision threshold.
% - Support-weighted F1 for the four PAM50 subtypes, whose class sizes run from 67 to 421 tumours.
% - MCC for the nucleotide sets. It is high only when all four entries of the confusion matrix are good \citep{Chicco2020}.
% - Definitions are in Chapter 2, "Evaluation methodology". No formulas here.
% - Metrics are never averaged across datasets or domains.
% Paragraph 5: how results are reported
% - Every score is on an outer test fold the search never saw. Report the per-fold scores and their spread.
% - Each hybrid is paired by fold with the best single representation on the same fold.
% - No significance tests and no confidence intervals, within or across datasets. With 5 folds, a p-value or an interval rests on 5 correlated values.
% - Fold scores share training molecules, so they are correlated. No unbiased estimator of the variance of k-fold cross-validation exists, and estimates that ignore the overlap underestimate it \citep{Bengio2004}.
% - Each comparison is reported as its per-fold differences and their spread across the 5 folds, with no threshold and no test.
% Paragraph 6: compute
% - Result needed: hours per dataset for the full search, and the hardware (CPU cores, memory, and a GPU where the DNN or the embeddings used one).
% ---- END NOTES ----

% ===== FROM: Comparators, controls and ablations notes =====
% ---- NOTES: Comparators, controls and ablations ----
% Takeaway: These experiments ask whether the hybrid's gain is real, or only the reward for searching many configurations.
% Takeaway: Where a search picks the best of many configurations, the comparator pays the same selection cost, by reselecting on the other folds.
% Paragraph 1: what these experiments are for
% - The hybrid picks the best of many configurations: allocation method, importance method, internal evaluator, budget.
% - A reader will ask whether the gain is the reward for searching hard. These experiments answer that.
% - Questions, one per paragraph below:
%   - does the search overfit; does a larger search help; fusing features or fusing predictions;
%   - the hybrid against standard feature reduction; must the internal evaluator match the final model;
%   - stability of the selection; gain against dataset size; importance during training;
%   - a few features added to one representation; a concatenation of two representations.
% - All use the group folds of "Datasets and splitting".
% Paragraph 2: the shared protocol, paying the same selection cost on both sides
% - The hybrid's score is the best of many configurations. Each single representation is one score.
% - The best of many cross-validated scores is biased upwards \citep{Varma2006}.
% - Remedy used in several experiments: leave-one-fold-out reselection.
% - For held-out fold $f$, pick the hybrid configuration with the best mean score on the other folds, $c^*_{-f} = \arg\max_c \frac{1}{F-1}\sum_{g\neq f} s_{c,g}$.
% - Pick the best single representation the same way, $k^*_{-f}$.
% - Report $\Delta_f = s_{c^*_{-f},f} - s_{k^*_{-f},f}$, signed so that $\Delta_f>0$ means the hybrid is better (flip the sign for RMSE).
% - It is arithmetic on stored per-fold scores. It needs no extra model fits.
% - Importance scores and the feature selection are always recomputed on each fold's training molecules only.
% - Some experiments choose the configuration by a search over all molecules and redo only the feature selection per fold. Their scores are not fully held out. Say so in each paragraph where it applies (marked below).
% Paragraph 3: does the search overfit? (held-out score against budget)
% - Nested cross-validation. For every budget $B$, that budget's winning configuration is rebuilt on the outer training molecules and scored on the outer held-out fold.
% - Selection bias at a budget = score the inner search reported minus the held-out score, signed so that a positive value means the search promised more than the held-out fold gave.
% - Budgets: the standard list (10 to 2,000 features), plus 4,000 and 8,000 features.
% - Also budgets at 0.5, 0.75, 1, 1.25, 1.5 and 2 times the smallest outer-fold training-set size (molecules). These place a budget where features equal molecules.
% - For the five downstream models the budget does not set the parameter count. The DNN and the random forest have more fitted parameters than training molecules at every budget. The SVM never can. The boosted trees' capacity is set by trees and depth.
% - So their curves say whether the search's budget choice overfits.
% - Two numbers per model: the budget each outer fold's search picked, and the budget with the best held-out score. A pick far above the best held-out budget means the search overfits.
% - Figure: x = budget / training-set size (log scale), y = held-out score, one line per model; second panel: selection bias against the same x.
% - Result needed: held-out score and selection bias against budget, per model and dataset.
% Paragraph 4: does a larger search give a better hybrid?
% - The size of a search is the number of configurations it scores: budgets x allocation methods x importance methods x internal evaluators x representation counts.
% - The question: does the hybrid's lead over each comparator (the ensemble, the best single representation, PCA, feature reduction at the same width) grow with the number of configurations?
% - The comparators get no search. Part of any gap that grows with search size is the hybrid taking the best of more draws.
% - Leave-one-fold-out reselection charges the hybrid for this. Report the reselected numbers.
% - In the full runs each dataset is searched at one tier only, full or fast. A search-size comparison needs the same dataset searched at two sizes.
% - Result needed: hybrid minus each comparator, per dataset, at each search size.
% Paragraph 5: fusing features against fusing predictions (ensembles)
% - Ensemble: one model per representation, predictions pooled. Hybrid: one model on one fused feature vector.
% - Joining features before one model is early fusion; pooling the predictions of one model per feature set is late fusion \citep{Stahlschmidt2022}.
% - The hybrid runs its own search on each outer fold's training molecules.
% - Every choice is made on an inner 3-fold group split of the training molecules. The held-out fold is used once.
% - Ensemble choices, all made on the inner split:
%   - the same model on every representation, or each representation with its own best model;
%   - members: the top 3 representations, or all of them;
%   - pooling: mean, median, trimmed mean, outlier-filtered means, score-weighted mean, or mean of the top 3.
% - Classification ensembles pool predicted probabilities.
% - Controls: the best single representation; all representations joined then reduced by PCA to the hybrid's width; the best (model, representation) pair read off the held-out fold. The last is an upper bound, not a fair comparator.
% - Cost, reported beside the score: representations computed at prediction time, models deployed, fit time and prediction time (seconds).
% - The hybrid only saves cost if it draws on fewer representations than the ensemble.
% - An ensemble is the obvious simpler alternative. If it matches the hybrid at lower cost, that belongs in Chapter 7, "Simple is often better".
% - QSAR comparison: \citet{Turkina2025} combine six non-hashed fingerprints by low-, mid- and high-level fusion on six regression and classification datasets. Mid-level fusion, which keeps fingerprint bits by their importance in the per-fingerprint models, gave the best accuracy.
% - TODO(lit): read the Turkina2025 methods for whether high-level fusion pools the predictions of one model per fingerprint, as the ensemble here does.
% - Result needed: held-out score and cost per method.
% Paragraph 6: the hybrid against standard feature reduction at the same width
% - Every method gets the same number of features $B$, drawn from all base representations, fitted on the fold's training molecules. Regression datasets only.
% - Methods: the hybrid (its stored configuration at budget $B$, features re-selected per fold); PCA; mRMR \citep{Peng2005}, whose ranking stops at 200 features for run time; recursive feature elimination with a 50-tree, depth-5 random forest \citep{Guyon2002}; $B$ random features.
% - PCA here: all base representations joined, each feature z-scored on the fold's training molecules, PCA fitted on those molecules. The number of components is $B$, capped at the total number of features and at one fewer than the number of training molecules. Held-out molecules are projected with the fitted PCA. The PCA control in Paragraph 5 is built the same way.
% - One downstream model for all, a 100-tree random forest, so only the features differ.
% - Not fully held out: the hybrid's configuration is chosen using all molecules. Only its feature selection is redone per fold.
% - It asks whether the gain comes from drawing on several representations, or only from having fewer features.
% - Figure: held-out score against $B$ (features), one line per method, one panel per dataset.
% - Result needed: held-out score against $B$ for each method.
% Paragraph 7: does the internal evaluator need to match the final model?
% - The search scores each representation with an internal evaluator to split the budget. The winner goes to a downstream model.
% - Nested cross-validation as in Paragraph 3. One search per outer fold per evaluator setting.
% - Settings, each against the matched evaluator:
%   - matched: same family as the downstream model, at the lighter evaluator settings;
%   - near family: random forest scored by LightGBM and the reverse, XGBoost by LightGBM, SVM by a Gaussian process, the DNN by a narrower DNN (64 and 32 units);
%   - far family: SVM for the tree models, LightGBM for the SVM, k-nearest neighbours for the DNN;
%   - the DNN scored by a wider DNN (512, 256 and 128 units);
%   - a fixed set of three evaluators, random forest, LightGBM and XGBoost, best one kept.
% - Report each setting minus matched, paired per outer fold, at the same number of features.
% - Evaluators with the downstream random forest's and LightGBM's exact settings are implemented but used by no search, so the full runs will not report them.
% - The answer sets whether one cheap evaluator is enough, or the search must be re-run for each model.
% - Result needed: each setting minus matched, per downstream model.
% Paragraph 8: stability of the selected features across folds
% - Selection on each of the $M$ folds gives a binary vector $\mathbf{z}_m \in \{0,1\}^d$. $d$ is the number of features in all base representations joined.
% - Nogueira stability \citep{Nogueira2018}: $\hat\Phi = 1 - \dfrac{\frac1d\sum_{f=1}^d s_f^2}{\frac{\bar k}{d}\big(1-\frac{\bar k}{d}\big)}$, with $s_f^2 = \frac{M}{M-1}\hat p_f(1-\hat p_f)$.
% - $\hat p_f$ is the fraction of folds that selected feature $f$. $\bar k$ is the mean number of features selected.
% - $\hat\Phi=1$: the same features every fold. $\hat\Phi \approx 0$: what random selection gives.
% - $d$ must be that full width, not the union of selected features.
% - Unstable feature sets would weaken any chemical reading of the blueprints.
% - Result needed: $\hat\Phi$ per dataset and model.
% Paragraph 9: learning curves, gain against dataset size
% - Subsamples of 10, 25, 50, 75 and 100% of the molecules, drawn as whole scaffold groups. Above 5,000 molecules, absolute sizes of 250 to 2,000 molecules instead.
% - 3 subsamples per size; the 100% point is drawn once, since every draw is the same. A full search on each, with budgets above the number of training molecules dropped.
% - Gain at each size = hybrid minus best single representation, with leave-one-fold-out reselection.
% - Small subsamples search fewer budgets. Part of any rise in gain with size may come from searching more configurations.
% - Figure: gain against number of molecules, one line per model, one panel per dataset.
% - Result needed: gain per dataset, model and size.
% Paragraph 10: importance used during training
% - The hybrid at 100 features, with random forest, LightGBM and XGBoost (the models that accept per-molecule weights).
% - Importance is computed on the training molecules with the importance method the search picked for that hybrid.
% - Five conditions:
%   - unweighted;
%   - features multiplied by the square root of importance, normalised so the weights average 1 (after standardising);
%   - molecules weighted by label agreement with their 10 nearest neighbours in the rescaled space;
%   - both together;
%   - only the features strictly above the median importance.
% - Molecule weight = 1 - |own label - neighbours' mean label| / label range, clipped to between 0.1 and 1.
% - Not fully held out: the configuration comes from a search over all molecules. Only the feature selection is redone per fold.
% - The molecule weighting is a noise-mitigation method: it down-weights suspect labels.
% - Tree ensembles are unchanged when a feature is multiplied by a positive constant: the splits depend only on the order of values \citep{Hastie2009}.
% - So the second condition and the unweighted one differ in one way only. Features with zero importance get weight zero, become constant, and drop out. Every other feature is only rescaled, which the three tree models ignore.
% - The rescaling does matter for the third condition: it changes which molecules are nearest neighbours.
% - Checking a label against its nearest neighbours goes back to editing a training set by its nearest neighbours \citep{Wilson1972}. Here the molecule is down-weighted, not removed.
% - No QSAR study was found that down-weights training molecules whose labels disagree with their nearest neighbours.
% - Nearest prior work, all outside QSAR. Nearest-neighbour noise filters adapted to regression labels remove such samples rather than down-weight them \citep{Martin2021}. Learned per-sample training weights for noisy labels come from general machine learning \citep{shu2019metaweightnet}.
% - Result needed: one result line.
% Paragraph 11: adding a few features to one strong representation, with decoys
% - Two decoys are planted among the candidate features: a block of pure Gaussian noise as wide as the real candidate features combined, and near-copies of 64 base features with noise at 5% of each feature's standard deviation.
% - Count how many decoy features are taken, per fold. A working rule takes zero.
% - Also run: the same extra features added to the multi-representation hybrid. It asks whether extra features do more on one representation or on the hybrid.
% - Result needed: one result line, with the decoy counts (features per fold).
% Paragraph 12: are hybrids more robust to label noise than their base representations? (only if this comparison stays in this chapter)
% - Robustness is Chapter 3's normalised area under the retention curve (\aucnorm, Equation~\ref{eq:auc_norm}).
% - Two differences from Chapter 3. First, the retention curve uses each dataset's own metric, not $R^2$; for an error metric the ratio is inverted, so 1 still means no loss.
% - Second, the noise levels: seven, from 0 to 1.5, each a fraction of the fold's clean training-label spread. Censoring has its own axis, the fraction of labels clipped (0 to 0.40).
% - The comparison: the hybrid at each budget (features) against each single representation and against the best single representation, re-picked in each fold. One downstream model, a random forest, under each regression noise type.
% - The hybrid's settings are searched again at every noise type and level, on the noisy training labels. Held-out labels stay clean.
% - Reference point: in Chapter 3 the model, not the representation, explained most of the variance in \aucnorm. Here the model is fixed, so only the representation changes.
% Paragraph 13: a concatenation of two representations against each alone
% - Question: does joining two whole representations score above the better one, or between the two?
% - For each dataset, take its top three single representations and form each of the three pairs. Score each representation alone and each concatenation, with the same model and folds.
% - A concatenation that lands between its two parts supports the case for combining and filtering at once.
% - Figure: one panel per dataset. For each pair: score of A, score of B, score of the concatenation, with the band between the two single scores shaded.
% - Result needed: for each pair, the score of each representation alone and of their concatenation.
% ---- END NOTES ----

% ===== FROM: Transfer across datasets, antibody mutational scans and data held at separate sites notes =====
% ---- NOTES: Transfer across datasets, antibody mutational scans and data held at separate sites ----
% Takeaway: Three setups test the search outside the single-dataset setting: reusing search settings on a second dataset, hybrids on antibody mutational scans, and feature selection with data held at separate sites.
% Takeaway: Each needs only several numeric representations of the same objects, one supervised target, and a grouped split.
% Paragraph 1: what this section covers
% - Three setups: reusing search settings on a second dataset, hybrids on antibody mutational scans, and feature selection with data held at separate sites.
% - These are method paragraphs: how each setup was built and scored, not what came out.
% - Each setup gets a short paragraph of background: what the task is and where it is used.
% - One requirement covers all three: several numeric representations of the same objects, one supervised target, and a grouped split that keeps near-duplicates on one side.
% Paragraph 2: reusing search settings across datasets, the question and the data
% - Question: can search settings found on a large dataset be reused on a small one, instead of running a new search?
% - The search settings are the importance method, the allocation method, the budget $B$ and the internal evaluator. They are decisions, not a set of features.
% - Source dataset: lipophilicity, 5,028 molecules. Target dataset: hERG binding affinity ($\mathrm{p}K_i$), 1,412 molecules.
% - Candidate representations are those listed for both datasets, kept only where the width (number of features) matches. Both searches then range over the same features.
% - The current candidate lists share seven representations: Avalon, ECFP4, Mordred and four pretrained embeddings. The width check can still drop some of them.
% - Result needed: the shared representation set the full run uses.
% - Why it matters: the search is the expensive step, and settings that transfer make it a one-off cost paid on the largest dataset.
% - Closest prior work: Meta-QSAR chose the learning method and molecular representation for a new QSAR dataset from results on over 2,700 other QSAR datasets. It beat random forest on a fingerprint by up to 13% on average \citep{Olier2018}. That reuses decisions, like the source search settings here, not a set of selected features.
% - No paper found so far applies feature indices selected on one QSAR dataset unchanged to another. The search was incomplete, so do not claim it has not been done.
% Paragraph 3: reusing search settings across datasets, the four comparisons
% - All four run inside every held-out fold of the target.
% - Best single representation, re-chosen on that fold's training molecules.
% - Hybrid searched on that fold's training molecules. This is the upper reference.
% - Source search settings, with features re-selected on the target's training molecules.
% - The source's own selected feature indices, applied unchanged. No target label is used to pick features.
% - The last two separate two questions: do the search decisions transfer, and do the chosen features transfer.
% - Five scaffold-group folds, RMSE in $\mathrm{p}K_i$ units, five downstream models (random forest, LightGBM, XGBoost, SVM, DNN).
% - Result needed: hybrid minus best single representation for each transfer comparison, per downstream model.
% Paragraph 4: hybrids on antibody mutational scans
% - Background, two sentences: a deep mutational scan measures a property (here expression or binding) for many single or few-point mutants of one antibody. It is used to map which positions tolerate change.
% - Four scans with a single wild type, including the expression scan among the datasets used here \citep{Koenig2017}.
% - Candidate representations: the protein language-model embeddings and cheaper sequence features of "Representations by domain", plus two representations of the change from the wild type. One is a 5-feature physicochemical difference; the other has the same five features at every sequence position.
% - Nested cross-validation: the search sees outer-training molecules only; the hybrid and every single representation are scored on the untouched outer fold.
% - Metric: Spearman correlation for the continuous scans, AUROC for the binarised binding scans.
% - Report hybrid minus best single representation, per scan, on the outer folds.
% - The two wild-type-difference representations are also scored as single representations, and each run reports their rank among all representations.
% - How many features each representation got is stored per fold with every result, but this comparison does not report it. So hybrid minus best single says nothing about the wild-type-difference representations unless that count is also reported.
% - Result needed: features the hybrid took from each of the two wild-type-difference representations, per scan, per fold.
% - Result needed: hybrid minus best single representation per scan, per downstream model.
% Paragraph 5: data held at separate sites, background
% - Federated learning trains one model across data held by several parties without moving the data \citep{McMahan2017}.
% - It has been run across pharmaceutical companies on shared property-prediction targets \citep{Heyndrickx2024}.
% - In this field the motivation is commercial: structures and assay results are the asset, so pooling them is not possible \citep{Hanser2023,Bassani2023}.
% - The collaborator (Lhasa Limited) uses the teacher-student process described in Chapter 2, "Noise in molecular data" \citep{Hanser2025}.
% - This thesis uses the collaborator's implementation unmodified, so the numbers stay comparable to theirs.
% Paragraph 6: how the teachers' predictions are combined
% - Distance weighting: each teacher's vote on a public molecule is weighted by how close that molecule is to the teacher's own training molecules, in fingerprint distance. A site counts most on chemistry it has seen.
% - Closeness: the mean Euclidean distance between fingerprints, from the query molecule to its 8 nearest neighbours among a random 20% of the site's training molecules.
% - These distances are rescaled to 0-1 across all sites, turned into a similarity as (1 - distance) squared, and normalised to sum to 1 over the sites.
% - TODO(lit): check that \citet{Hanser2025} states this weighting in its methods before citing it for it; otherwise cite the collaborator's released code.
% - The student is trained on the most confident combined labels, the same number from each class.
% - Score weighting: each site gets one weight from its own cross-validated score, in place of the per-molecule distance.
% - Only predictions on shared public molecules cross a site boundary.
% Paragraph 7: feature selection at the sites
% - The addition in this thesis: feature selection, not only model fitting, happens at the sites.
% - Each site runs the hybrid search on its own molecules. Features are chosen on 80% of its molecules and scored on the other 20%, split by scaffold group.
% - Each site reports three things: feature indices, importance values and one score. No molecule and no feature value leaves the site.
% - A coordinator splits the total budget $B$ across sites in proportion to their scores. Each share $b_s$ is rounded down, with at least one feature per site, so $\sum_s b_s \approx B$.
% - Only the proportional split is run. With no equal-split comparison, the text cannot say what the site scores buy.
% - A feature another site has already taken is skipped and the site takes its next-ranked feature instead, so the features in the hybrid are all distinct.
% - Why it matters: indices and importances are smaller, and harder to invert, than model weights or gradients.
% - Why it matters: splitting the budget between sites is the same question the allocation method answers between representations.
% - Related work: each client runs a mutual-information feature selection on its own data and a server aggregates the results, with no raw data exchanged \citep{Cassara2022}.
% Paragraph 8: the hERG comparison, data
% - Background, one sentence: hERG channel blockade is a cardiotoxicity liability screened early in drug discovery, and its public data comes from many labs \citep{Czodrowski2013}.
% - Three real sources act as sites: the public ChEMBL hERG $K_i$ set used elsewhere in this thesis, and the collaborator's training set (7,772 molecules) and test set (4,383 molecules).
% - The two collaborator sets were labelled by different curation rules, so the sites differ in curation, not only chemistry.
% - This is source heterogeneity that exists in the world. The process is usually simulated with sites made by clustering one dataset.
% - Molecules held by more than one source are deduplicated on InChIKey, keeping the public copy.
% - Binary target for the public set: blocker if $\mathrm{p}K_i \ge 6$, i.e. $K_i \le 1\,\mu\mathrm{M}$.
% - Five scaffold-group folds over the pooled molecules.
% - The search is scored by AUROC. The reported number is MCC, which does not reward predicting the majority class.
% Paragraph 9: the hERG comparison, the two-by-two and the context comparisons
% - Molecules pooled or kept at their sites, crossed with all features or a selected set.
% - Pooled molecules, all features: the upper reference, not possible in practice.
% - Pooled molecules, selected features: what the selection is worth on its own.
% - Molecules at their sites, all features: federated training as usually done.
% - Molecules at their sites, selected features: the only setting where nothing but indices, importances, one score and predictions moves.
% - Read the corners: selected minus all features, at fixed site structure, is what selection buys. Pooled minus site-held, at fixed features, is what pooling buys.
% - Context comparisons: the single best source picked on training molecules; both teacher-weighting schemes; the same search on one fingerprint alone; one teacher per representation.
% - The budget $B$ is swept (50, 100, 250, 500 and 1,000 features), so no comparison rests on one width.
% - Caution: both all-features settings fit one 2,048-bit Morgan fingerprint (radius 2, pharmacophoric atom types). Both selected settings fit features drawn from the hERG $K_i$ dataset's candidate representations.
% - So selected minus all features differs in two ways at once, width and source. The budget sweep varies the width of the selected set only. It cannot separate the source difference.
% - The single-fingerprint search is the closest control for the source difference: same search, ECFP4 only.
% - Two curators disagreeing on one molecule is label noise with a known source (Chapter 6).
% - Figure: two rows, the same three sites in both. Top: sites, teacher models, a shared public set being labelled, one student. Bottom: sites, each site's feature search, the coordinator splitting the budget ($b_1,b_2,b_3$ with $\sum_s b_s \approx B$), each site's model on the shared feature list, distance-weighted predictions. Mark what crosses a site boundary (top: predictions on public molecules; bottom: feature indices, importances, one score, predictions).
% - Figure: the hERG two-by-two. Horizontal axis molecules pooled or site-held, vertical axis all or selected features, one box per setting, arrows marked as the selection difference and the pooling difference.
% Paragraph 10: other separate-site experiments, one line each
% - Sites with different amounts of label noise: selection at the sites against the teacher ensembles and one model trained on all sites' molecules pooled, at several noise levels. One version has each site remove the molecules it flags as noisy first.
% - Representations as sites: every site holds the same molecules as a different representation. Feature-level fusion is compared with pooling the teachers' predictions.
% - Sites holding different families of representations (substructure, physicochemical, pretrained) of shared molecules.
% - Antibody sources as sites: a binding-energy-change set, a deep mutational scan of binding, and a thermostability set. Within a site, sub-sites are made by clustering language-model embeddings, not fingerprints.
% - In the antibody version, the two teacher-ensemble comparisons are skipped. Each site's features are reduced separately, so one site's teacher cannot predict on another site's features. The labels are also in incompatible units (kcal/mol, a normalised score, degrees Celsius).
% - Result needed: one line on the outcome of each.
% Paragraph 11: limits of the separate-site work
% - One source is proprietary and cannot be redistributed. The comparison cannot be reproduced from public data alone (Appendix~\ref{app:availability}).
% - The comparison depends on the collaborator's implementation of the process.
% - The published process labels a large separate public set. Here the training molecules play that role, and held-out molecules are excluded from it. State this as a compromise.
% - Sites differ in size and in how their labels were derived, so a site's score reflects its curation as well as its chemistry.
% Paragraph 12: possible later additions (external benchmark tasks, not among the current datasets; not now)
% - Activity-cliff prediction on dopamine D2 receptor, factor Xa and SARS-CoV-2 main protease, following a published comparison of representations and models \citep{Dablander2023}.
% - Antibody developability on the biophysical properties of a public competition set, using that competition's cluster splits \citep{vanNiekerk2026}.
% - Antibody--antigen binding free-energy change: how close sequence representations get to a structure-based graph network \citep{Hummer2025}.
% ---- END NOTES ----

% ===== FROM: Results: when the hybrid helps notes =====
% ---- NOTES: Results: when the hybrid helps ----
% Takeaway: Dataset by dataset, where the hybrid scores better than the best single representation and than an ensemble, and where it does not.
% Takeaway: Where it does not help, one representation may already carry most of the signal; a tie there is a sign the search is not inflating scores.
% Paragraph 1: how to read every number in this section
% - Every comparison is on the outer held-out folds of the group splits.
% - Report the leave-one-fold-out reselected numbers for both the hybrid and the best single representation. Reselection charges both sides for being a maximum.
% - Each dataset is reported in its own metric. Give per-fold scores and their spread. No significance tests, and no averaging across datasets.
% Paragraph 2: hybrid against the best single representation, dataset by dataset
% - One row per dataset and downstream model: the best single representation and its score, the hybrid's score, and hybrid minus best single, fold by fold.
% - Sign every difference so that positive means the hybrid is better, in both metric directions.
% - Group the rows by the six domains: small molecules, the quantum-chemical property, reactions, antibodies, nucleotide sequences, cancer genomics.
% - Report the count of datasets, per model, where the reselected hybrid is ahead on the fold mean. Say what the top and bottom of the count are (datasets ahead, datasets run).
% - The best single representation is chosen per fold, so it can differ between folds. Name it in each fold.
% - Table (\label{tab:hybrid_vs_single}): one row is one dataset scored with one downstream model. Columns: domain; dataset; metric; the best single representation in each fold; its score in each fold; the hybrid's score in each fold; the difference in each fold (positive = hybrid better); mean and spread of the difference across the 5 folds.
% - Result needed: the per-dataset table from the full runs.
% Paragraph 3: hybrid against an ensemble of one model per representation
% - Report three differences per (dataset, model): hybrid minus the ensemble that chooses its own members and pooling rule on the inner split, hybrid minus best single, and the upper bound read off the held-out fold.
% - The upper bound is not a fair comparator. Report it only to show how far the fair comparators are from it.
% - Report the cost beside the score: representations computed at prediction time and models deployed. The hybrid saves only where it draws on fewer representations than the ensemble.
% - If an ensemble matches the hybrid at lower cost on a dataset, say so here; it also belongs in Chapter 7, "Simple is often better".
% - Result needed: hybrid, ensemble and best single per dataset and model, with the cost figures.
% Paragraph 4: how the gain changes with the number of molecules
% - Report hybrid minus best single at each size, reselected, per dataset and model.
% - At the smallest sizes the training folds can hold fewer molecules than the smallest budget. Those sizes search fewer budgets. Say at which sizes this happens, per dataset.
% - So a gain that rises with size can come from more data or from searching more budgets. The curve alone cannot separate the two. State this next to the figure.
% - Figure: x = number of molecules (log scale), y = hybrid minus best single in the dataset's metric, positive = hybrid better; one line per model; one panel per dataset; per-fold points shown.
% - Result needed: gain against number of molecules, per dataset and model.
% Paragraph 5: does a larger search give a larger gain? (only if paired runs at two search sizes are made)
% - A search-size comparison needs the same dataset searched at two numbers of configurations, with folds and representations unchanged. The full runs search each dataset at one tier only.
% - An early comparison of a larger search (5 folds, 16 configurations per model) against a smaller one (3 folds, 2 to 6 configurations) is not usable. It differs in three ways at once, not one: the number of configurations, the number of folds, and the representations searched.
% - If the paired runs are made: report hybrid minus each comparator at the two sizes, reselected, per dataset and model. The two tiers differ in which allocation methods, importance methods and representation counts are searched, not only in how many configurations. Say so next to the numbers.
% Paragraph 6: where the hybrid does not help
% - Name the datasets and models where the reselected hybrid is behind or level with the best single representation.
% - For each, report what share of the hybrid's features comes from the best single representation. That count is stored per fold with every searched configuration, so it needs no extra run.
% - This tests one explanation: where one representation carries most of the signal, the search can only add features from weaker ones. Report the share; do not state the explanation as found.
% - Also report training-fold size and the smallest budget searched, since small datasets search fewer budgets.
% - Result needed: the datasets where the hybrid is not ahead, with the best representation's share of the hybrid's features.
% ---- END NOTES ----

% ===== FROM: Results: why the hybrid helps notes =====
% ---- NOTES: Results: why the hybrid helps ----
% Takeaway: The gain comes from which features are chosen, not only from having fewer: at the same number of features the hybrid is ahead of PCA and random features.
% Takeaway: Different downstream models want different numbers of features and different representations.
% Paragraph 1: what this section asks
% - Five questions, one paragraph each. Does the gain come from drawing on several representations, or only from having fewer features? Must the internal evaluator be the downstream model? Does importance help when used in training? Is adding a few features to one strong representation safe? Do different downstream models want different hybrids?
% - Scores are on held-out group folds, in each dataset's own metric. Give per-fold scores and their spread. No significance tests.
% Paragraph 2: several representations, or only fewer features?
% - If the hybrid leads PCA and random features at the same $B$, the gain is not only from having fewer features.
% - Report the realised number of features for every method at each $B$. Compare only where the widths match. At large $B$ the hybrid can end with fewer than $B$ features.
% - One sentence of caveat: the hybrid's configuration was chosen using all molecules.
% - Preliminary (to be replaced by the full runs), widths matched, PCA, recursive feature elimination and random features only (no mRMR). Antibody thermostability (Tresanco), Spearman, hybrid / PCA / recursive feature elimination / random features:
%   - 10 features: 0.5945 / 0.4860 / 0.4682 / 0.2145.
%   - 50 features: 0.6131 / 0.5225 / 0.5689 / 0.5094.
%   - 250 features: 0.6029 / 0.4824 / 0.6004 / 0.5636.
% - Same runs, hERG Ki at 10 features: hybrid 0.6304 RMSE (pKi) against 0.6776 for PCA; the hybrid scored better on all 5 folds.
% - Same runs: at 1,000 features the methods ended at different widths, so that budget cannot be compared.
% - Figure: x = $B$ (features, log scale), y = held-out score in the dataset's metric, one line per method, one panel per dataset.
% - Result needed: held-out score against $B$ for each method, per dataset.
% Paragraph 3: must the internal evaluator be the downstream model?
% - Report each evaluator setting minus the matched evaluator, paired per outer fold, at the same number of features.
% - For random forest, LightGBM and XGBoost the three-evaluator set contains the matched evaluator. It cannot score worse on the inner folds than matching.
% - The finding to state, whichever way it goes: which evaluator best predicts what the downstream model will do with a feature set.
% - If the matched evaluator is not always best, read it as informative, not as a failure. The default search tries four internal evaluators and keeps whichever scores best, so it already covers both outcomes.
% - Why the same number of features matters: a difference between settings that ended at different widths is not a difference in evaluator.
% - No numbers exist yet for near family, far family or the wider network.
% - Table: one row per downstream model, per dataset. Columns, each minus matched (mean and spread over folds): near family, far family, wider network (DNN row only), three-evaluator set.
% - Result needed: each setting minus matched, per downstream model and dataset.
% Paragraph 4: does importance help when used during training?
% - If the feature-scaling condition stays, read it as "delete the zero-importance features" on these three tree models.
% - Report how many of the hybrid's features get zero importance, per dataset (features out of 100). This says how much the condition can change.
% - The same scaling used as a distance for nearest-neighbour prediction is reported in Chapter 6, not here.
% - Preliminary, 250-feature hybrid: LightGBM gain gave exactly zero importance to 66 to 78 of the 250 features.
% - Same runs, Caco-2: every weighting tried scored at or below unweighted, for random forest and for LightGBM.
% - Result needed: each condition minus unweighted, per model and dataset, and the zero-importance counts.
% Paragraph 5: adding a few features to one strong representation
% - Base: each of the dataset's ten representations in turn, used whole. Report the best single representation per dataset in the main table.
% - Control at the same width: the same number of extra features drawn at random from the other representations. With-extras minus random-extras separates the selection from simply adding features.
% - Report the score with and without extras, per model and dataset.
% - Report how many of the features taken per fold are real, how many are near-copies of base features, and how many are pure noise.
% - If decoys make up a large share of what is taken, say that a score gain alone does not show the extras carry signal.
% - Report whether extras do more on one representation or on the multi-representation hybrid.
% - Preliminary, antibody thermostability (Tresanco, 672 antibodies), 3 folds, Spearman. Base: SaProt alone (1,000 features). Candidate extras: the 48 engineered mutation features. Up to 20 extras taken; no admission rule.
%   - Random forest: 0.5683 without extras, 0.5762 with them.
%   - SVM: 0.5716 without extras, 0.5830 with them.
%   - Features taken per fold, random forest, 3 folds: real 12, 7, 18; near-copy decoys 21, 11, 21; noise decoys 7, 2, 1.
% - Table: one row per (dataset, model). Columns: base representation, score without extras, score with extras, real features taken, near-copy decoys taken, noise decoys taken (per fold).
% - Result needed: score with and without extras, and decoy counts (features per fold), per dataset and model.
% Paragraph 6: do tree ensembles and neural networks want different hybrids?
% - Two parts: the number of features each model does best with, and which representations each model does best with.
% - Number of features: for each downstream model, the budget with the best held-out score.
% - Representations: each whole base representation is fed to each downstream model, 5 folds. Rank the representations by score for each model. These scores come from the full runs.
% - Report the Spearman correlation of those rankings between each pair of models, per dataset (about 10 representations per dataset).
% - A low correlation between a tree model and the DNN means the two models would fund different representations. Name the best representation per model where they differ.
% - Then compare the features chosen by a tree-model hybrid and a DNN hybrid on the same dataset: overlap as Jaccard, $|A\cap B|/|A\cup B|$, over the selected features, broken down by source representation.
% - The search is run separately for each downstream model, so each model gets its own budget and representations.
% - Preliminary, nanobody thermostability, 5 folds: SVM did best at 300 features (0.6384 Spearman); DNN, LightGBM and random forest at 50 features; XGBoost at 100.
% - Figure: heatmap of rank correlation between downstream models, one panel per dataset.
% - Figure: for one dataset, held-out score against budget (features, log scale), one line per downstream model, showing where each peaks.
% - Result needed: best budget per model and dataset; rank correlation per model pair and dataset; Jaccard overlap per dataset.
% ---- END NOTES ----

% ===== FROM: Results: blueprints on specific datasets notes =====
% ---- NOTES: Results: blueprints on specific datasets ----
% Takeaway: On datasets with named features, the best hybrid's selected features are read against the known chemistry of the endpoint.
% Takeaway: Which settings win differs by dataset, and the selection is only partly stable across folds.
% Paragraph 1: what this section does
% - Take the best hybrid on a few datasets and read its selected features against the known chemistry of the endpoint.
% - Only named features can be read chemically. ECFP bits, MACCS keys and the physicochemical descriptor vector (PDV) are named; Mordred, atom-pair and Avalon features are indices only.
% - So choose datasets whose representations include ECFP4 or PDV: hERG Ki, Caco-2 permeability or lipophilicity.
% - Of each dataset's 10 representations, hERG Ki and Caco-2 carry two readable ones (ECFP4 and PDV). Lipophilicity carries one (ECFP4).
% Paragraph 2: how the features are read
% - Two downstream models are used: random forest and XGBoost.
% - Per fold, each of the two is trained on its own hybrid, built from its own best search result. No model's hybrid is read under another model.
% - SHAP values are computed only for the model that scores best on that fold, on the fold's held-out molecules. Each feature's score is its mean absolute SHAP value.
% - Both models are tree ensembles, so SHAP values come from TreeSHAP \citep{Lundberg2020}.
% - Every selected feature is mapped back to its representation. Absolute SHAP values are summed per representation and scaled to sum to 1: each representation's share of the hybrid.
% - The shares are computed inside each fold, then averaged over folds. The list of top features comes from one fold, because each fold selects different features.
% - An ECFP bit is resolved from the first of the dataset's first 50 molecules that sets it, and that molecule's atom environment is drawn. A bit none of those 50 molecules sets stays unresolved.
% - One hashed bit can stand for several environments, so say which one was drawn.
% Paragraph 3: one worked example, hERG Ki
% - The budget split across representations: how many features each representation contributes, and its share of summed importance.
% - The top named features, with their substructures or descriptor names.
% - Read them against what is known about hERG binding: most blockers carry a basic nitrogen, protonated at physiological pH, and hydrophobic or aromatic groups. In the channel pore these meet two aromatic residues, Tyr652 and Phe656.
% - Sources, neither yet in refs.bib: Aronov, "Predictive in silico modeling for hERG channel blockers", Drug Discovery Today 10(2):149--155, 2005, doi 10.1016/S1359-6446(04)03278-7 (hERG blocker pharmacophore models). Mitcheson et al., "A structural basis for drug-induced long QT syndrome", PNAS 97(22):12329--12333, 2000, doi 10.1073/pnas.210244497 (the Tyr652 and Phe656 binding residues).
% - State for each top feature whether it confirms known chemistry, suggests something not expected, or cannot be read (unnamed feature).
% - Selected means predictive, not causal; correlated features can stand in for one another. One sentence.
% - Result needed: per-representation feature counts and importance shares for the hERG Ki hybrid, for the model explained on each fold.
% - Result needed: the top 10 to 20 named features with their SHAP values.
% Paragraph 4: a second dataset, for contrast
% - Same reading as Paragraph 3, shorter: does the selection pick features known to drive permeability or lipophilicity?
% - Result needed: as for hERG Ki.
% Paragraph 5: how stable the selection is across folds
% - Measured with the Nogueira stability index. 1 means the same features in every fold; about 0 is what random selection gives.
% - Preliminary, 3 folds: 0.447 on Caco-2 permeability, 0.275 on hERG Ki, 0.074 on nanobody thermostability (Tresanco).
% - Each of these runs used 5 of the dataset's 10 representations, and only random forest and SVM as downstream models. For thermostability all 5 were learned embeddings. The numbers will change with the full set.
% - Preliminary: the search was also repeated with one more fold held out, and the result scored on that fold. The score was worse on every fold of all three datasets.
% - Caco-2: RMSE 0.2718 to 0.2953, 0.2986 to 0.3289, 0.2740 to 0.3077 (own fold to one fold fewer). hERG Ki: RMSE in pKi 0.5702 to 0.6392, 0.5862 to 0.6636, 0.5605 to 0.6034.
% - Thermostability: Spearman 0.6335 to 0.5743, 0.6703 to 0.6529, 0.5530 to 0.5286.
% - This comparison changes two things at once: the search sees one fold fewer of molecules, and it produces different features. It does not test whether features chosen on one fold carry over to another fold.
% - What a low index means for the chemical reading: correlated features can swap between folds and carry the same information.
% - Per-fold spread only; no confidence intervals and no significance tests.
% - Table: stability index, one row per dataset, one column per downstream model.
% - Result needed: stability index per dataset and downstream model, 5 folds.
% Paragraph 6: does the blueprint change across chemical space?
% - Molecules are split into 5 k-means clusters of chemical space, fitted on the training molecules only. Importance and allocation are computed separately in each cluster.
% - Question: do different regions give each representation a different share of the budget? For example, do substructure fingerprints take more of the budget in regions with many activity cliffs?
% - Compare the per-representation budget shares between clusters with a Spearman rank correlation, averaged over cluster pairs.
% - Control: the same training molecules split at random into parts of the same sizes. Clusters matter only if their shares agree less than the random parts do.
% - Cluster edges: the clustering is repeated with a fraction of the training molecules (0, 0.01, 0.05, 0.1 and 0.2) moved at random to another cluster. Say whether the per-cluster shares change as the edges blur.
% - Also compare the selected features themselves between clusters: two clusters can split the budget the same way and still take different features.
% - The same question on scaffold-defined subsets of a dataset: the best search settings found in each subset.
% - Result needed: per-cluster budget shares by representation, hERG blockers (Lhasa Limited), against the random-partition control.
% Paragraph 7: do the best search settings co-vary across datasets?
% - For each dataset and downstream model: the winning allocation method, importance method, internal evaluator, representation count and budget.
% - Which datasets land on the same settings, and whether data complexity measures computed on each dataset track the winning budget. One Spearman coefficient per complexity measure and downstream model; coefficients only, no p-values.
% - This is the one place the full-run winners per dataset are reported.
% - Figure: one small panel per setting (allocation method, representation count, importance method, internal evaluator). One bar group per dataset; bar = how many (downstream model, budget) combinations each value won.
% - Figure: for one hERG Ki hybrid. Left: stacked bar of the budget split across representations (features per representation). Right: the top named features, with ECFP and MACCS substructures drawn on an example molecule and descriptor names listed.
% - Table: one row per dataset (hERG Ki and the contrast dataset). Columns: downstream model, budget in features, features per representation, top 5 named features.
% - Result needed: table of winning settings, one row per dataset, one column group per downstream model.
% ---- END NOTES ----

% ===== FROM: Results: transfer across datasets and antibody mutational scans notes =====
% ---- NOTES: Results: transfer across datasets and antibody mutational scans ----
% Takeaway: Search settings found on lipophilicity, with features re-selected on hERG, beat hERG's best single representation on 3 of 5 downstream models; copying the features unchanged did so on 2 of 5 (preliminary).
% Takeaway: On an antibody expression scan the hybrid was ahead of the best single representation on all 5 downstream models (preliminary).
% Paragraph 1: what this section reports
% - First: search settings found on lipophilicity reused on hERG $\mathrm{p}K_i$. Second: the hybrid against the best single representation on an antibody deep mutational scan.
% - All numbers below are preliminary and will be replaced by the full runs.
% Paragraph 2: reusing search settings, the numbers
% - Metric: RMSE in $\mathrm{p}K_i$ units on hERG, lower is better. Five scaffold-group folds.
% - Each difference is the best single representation's RMSE minus the comparison's RMSE, mean over the 5 folds. Positive means the comparison had the lower error.
% - The best single representation is chosen per fold. It was not the same representation on every fold and model.
% - Re-selecting features on hERG with the source search settings beat the best single representation on 3 of 5 downstream models. Example: SVM, 0.0938 lower RMSE.
% - Copying the source's selected features unchanged beat it on 2 of 5 models. Example: SVM, 0.0791 lower RMSE.
% - Against the hybrid searched on hERG itself: re-selecting with the source settings came within 0.01 RMSE of it on 4 of 5 models (SVM 0.5678 against 0.5693). The gap is 0.0121 on DNN, where the source settings scored lower.
% - Give the spread with each number. The fold standard deviation of the copied-features difference on DNN is 0.1823 RMSE, larger than any mean difference in the table.
% - State what the two comparisons separate: whether the search decisions carry over, and whether the chosen features carry over. Let the table carry the numbers; no verdict beyond the counts.
% - Table: one row per downstream model (random forest, XGBoost, SVM, DNN, LightGBM). Columns: best single representation RMSE; hybrid searched on hERG RMSE; source settings re-selected RMSE; source features copied RMSE; difference for each of the last two against the best single, with fold standard deviation. Caption: metric, units, 5 folds, positive means lower error than the best single representation.
%   Source settings, features re-selected: SVM +0.0938, DNN +0.0651, LightGBM +0.0300, random forest -0.0037, XGBoost -0.0222.
%   Source features copied unchanged: SVM +0.0791, LightGBM +0.0314, DNN -0.0222, random forest -0.0203, XGBoost -0.0052.
% - Result needed: the same table from the full runs.
% Paragraph 3: the antibody mutational scan, the numbers
% - Dataset: the expression scan \citep{Koenig2017}, 4,275 antibody variants of one wild type, on average 1.0 mutation per variant.
% - Nested cross-validation with 3 outer folds of unequal size: 874, 1,710 and 1,691 test variants (3,401, 2,565 and 2,584 training). Reduced search: 200 search configurations, budgets of 100 and 500 features.
% - Metric: Spearman correlation, higher is better. Difference is hybrid minus best single representation, mean over the 3 outer folds.
% - The hybrid was higher than the best single representation on all 5 downstream models. Example: LightGBM, +0.0431 Spearman.
% - Per outer fold, the hybrid was higher in 14 of 15 model-fold pairs (5 models, 3 outer folds). The exception is XGBoost on one fold (-0.030 Spearman).
% - The best single representation was ProtT5 on every model and every fold. Its held-out Spearman, mean over the 3 outer folds, ranged from 0.6511 (XGBoost) to 0.7154 (SVM).
% - The full run covers all four scans with the full search and the full list of budgets. The binding scans are scored by AUROC, not Spearman.
% - Table: one row per downstream model. Columns: best single representation (name, Spearman); hybrid Spearman; difference, mean and the 3 per-fold values. Caption: 3 outer folds, nested cross-validation, reduced search.
%   DNN +0.0286, LightGBM +0.0431, random forest +0.0366, SVM +0.0287, XGBoost +0.0257 (hybrid minus ProtT5, Spearman).
% - Result needed: the other three scans, and the full search on the expression scan.
% Paragraph 4: which features the antibody hybrids selected
% - Averaged over the 15 selected hybrids (5 models, 3 outer folds), the budget went to ProtT5 (33% of features), IgT5 (29%), AbLang-2 (24%) and ESM-2 (12%).
% - One-hot and BLOSUM sequence encodings together took 0.9% of features.
% - The two representations of the change from the wild type were selected by 6 of the 15 hybrids. Where selected, the two together took at most 6% of the features.
% - So report this as a hybrid-against-best-single result on antibodies, not as a result about the wild-type-difference representations. One line on them at most.
% - Figure (optional): stacked bar of the budget split across representations for each of the 15 antibody hybrids.
% ---- END NOTES ----

% ===== FROM: Results: data held at separate sites notes =====
% ---- NOTES: Results: data held at separate sites ----
% Takeaway: Does feature selection run separately at each site match feature selection on pooled molecules, when only feature indices, importances, one score per site and predictions cross between sites?
% Paragraph 1: how to read this section
% - The main comparison uses three real hERG sources as sites. The other separate-site experiments get one line each.
% - Every number is MCC on the held-out folds (five scaffold-group folds). Give per-fold scores and their spread. No significance tests.
% - The budget is swept at 50, 100, 250, 500 and 1,000 features. Report every budget, not one.
% Paragraph 2: the hERG two-by-two
% - The four settings: molecules pooled or kept at their sites, crossed with all features or a selected set.
% - Selected minus all features, at fixed site structure: what the selection buys.
% - Pooled minus site-held, at fixed features: what pooling buys. The pooled settings need all molecules in one place, so they are references, not options open to the sites.
% - Report the selection difference separately for pooled and for site-held molecules. Say at which budgets it keeps its sign.
% - Table (\label{tab:sites_two_by_two}): one row per budget (features). Columns: pooled/all, pooled/selected, site-held/all, site-held/selected, selected minus all (pooled), selected minus all (site-held). Mean MCC and spread over the five folds.
% - Figure: the two-by-two, with the MCC of each setting at one budget written in its box. Per-budget numbers stay in the table.
% - Result needed: MCC for the four settings at each budget, per fold.
% Paragraph 3: how to read the selection difference
% - Selected minus all features differs in two ways at once: the number of features and the representations they come from.
% - With pooled molecules, pooled/selected minus the ECFP4-only search, at equal width, is what drawing on several representations buys.
% - With site-held molecules, report the selection difference as a mix of width and source.
% - Result needed: how many representations reach all three sources (6 of 10, preliminary).
% - Result needed: pooled/selected minus the ECFP4-only search, at each budget.
% Paragraph 4: the context comparisons
% - Best single source, picked on its own training molecules and scored on the fold. Its gap to the site-held selected setting is what combining sites is worth.
% - Also report the best single source picked on the held-out fold. It needs the answers, so it is a bound, not a comparator.
% - The collaborator's two teacher-weighting schemes: per-molecule fingerprint distance, and one weight per site from its own cross-validated score.
% - One teacher per representation, predictions weighted by each representation's cross-validated MCC.
% - How the coordinator split the budget: each site's share is in proportion to the score it reported. Report each site's score and its share of the features, per budget.
% - A naive pool (one model on all molecules and the fingerprint) is the pooled/all corner. Do not report it twice.
% - Result needed: MCC of each context comparison at each budget; each site's score and share of the features.
% Paragraph 5: other separate-site experiments, one line each
% - Include those that came out null; say what was measured and what the difference was, in the experiment's own metric.
% - Sites with different amounts of label noise: selection at the sites against the teacher ensembles and a naive pool, at several noise levels. One version has each site remove the molecules it flags as noisy first.
% - Representations as sites: every site holds the same molecules as a different representation. Feature-level fusion against pooling the teachers' predictions.
% - Families of representations as sites (substructure, physicochemical, pretrained): one selection over all families against a selection per family, at the same total budget.
% - Antibody sources as sites: a binding-energy-change set, a deep mutational scan of binding, and a thermostability set. Report per-site scores, and the transfer of one site's selected feature positions to another.
% - The antibody sites' labels are in different units (kcal/mol, a normalised score, degrees Celsius). Two teacher-ensemble comparisons are not defined there and were not run. Say so in one sentence.
% - Copying a feature position between sites assumes the feature means the same at every site. State the antibody transfer result only for representations where that holds.
% - Result needed: one number per experiment, in its own metric, at each budget.
% Paragraph 6: what these results can and cannot support
% - Sites differ in curation as well as chemistry, so a gap between sites is not only a chemistry effect.
% ---- END NOTES ----

% ===== FROM: Discussion notes =====
% ---- NOTES: Discussion ----
% Takeaway: Feature selection can overfit, and the search picks the best of many configurations. Nested cross-validation charges the hybrid for that, so the held-out folds, not the inner search, decide whether it helps.
% Takeaway: Where the hybrid does not help, a tie with the best single representation is a sign the search is not inflating scores.
% Paragraph 1: selection bias is a known risk of feature selection (her paragraph below)
% - Keep the argument. Selection done outside the folds inflates the score \citep{Ambroise2002}.
% - Selection is a form of model selection with its own cost \citep{Cawley2010}. Searching adds degrees of freedom \citep{Tibshirani2015}. Stability can be measured \citep{Nogueira2018}.
% - Change: Ambroise and McLachlan's no-signal test was colon data (62 tissues), class labels randomly permuted, SVM with recursive feature elimination. With selection outside the folds, cross-validated error was close to zero at 128 genes and about 10\% at 8 genes. With selection inside each fold, the error estimates were 0.40 to 0.45. The "median error was 11 out of 40" figure does not match this.
% - Change: Vabalas et al. say selection on pooled training and test data adds "considerably more" bias than parameter tuning \citep{Vabalas2019}.
% - Change: write "The hybrid search", not "KIRBy". The search does not range over arbitrary feature subsets; it scores only the configurations in the grid, so "many possible feature subsets" should read "many configurations".
% - Change: the \citet{Ng2004} sentence misstates the paper. For logistic regression with $L_1$ regularisation, the number of training examples needed grows only with the logarithm of the number of irrelevant features. For rotationally invariant methods ($L_2$-regularised logistic regression, SVMs, neural networks) it grows at least linearly in the worst case.
% - Change: \citet{Kumar2022} found fine-tuning on average 7\% less accurate out of distribution than a linear probe on frozen features; "up to 10\%" should be checked against that.
% Paragraph 2: the case that the search overfits
% - Features are ranked by importance computed on the same training labels the model is then fitted to. Noise in those labels can steer which features are picked.
% - The search picks the best of many configurations. The best of many noisy scores is biased upwards, even when each score is unbiased \citep{Cawley2010}.
% - Prior work: Sort \& Slice, which uses no labels to choose fingerprint substructures, beat two supervised selection methods, filtering and mutual-information maximisation \citep{sns}.
% - The authors say the reasons are "not obvious". They call it "conceivable" that using the label to select substructures adds to the risk of overfitting. A hedged suggestion, not a measured finding.
% - Small datasets are the most exposed. With few molecules, one fold's importance ranking rests on few labels.
% Paragraph 3: the case that it does not, or that the cost is outweighed
% - The held-out test is nested cross-validation. The search runs on the outer training molecules only; its choice is scored on an outer fold it never saw.
% - The single representations are scored on the same outer folds.
% - Picking the best single representation out of $K$ is also a selection step with its own optimism. The hybrid compares far more configurations, so its selection cost is larger. Nested cross-validation charges each side for its own.
% - Stability across folds is a second check. Correlated features can swap between folds while carrying the same information, so a low index alone does not show overfitting.
% - Reusing search settings on a second dataset is a third check. Settings chosen on one dataset keep their score on another only if the selection was not fitted to noise.
% - Figure: x = budget $B$ (features, log scale), vertical line at the number of training molecules. y = score in the dataset's own metric. Two lines: the score the inner search reported, and the held-out outer-fold score. One panel per dataset. Individual fold scores as points; no error bars.
% - Result needed: selection bias per dataset and budget; whether the held-out score peaks and falls once $B$ exceeds the number of training molecules; the stability index per dataset.
% Paragraph 4: robustness from drawing on several representations
% - A hybrid draws its features from several representations. The model does not depend on one representation's blind spots, such as hash collisions or molecules it cannot tell apart (Chapter 4).
% - The tradeoff: part of any chosen set of $B$ features may be picked on noise. The question is whether strong, varied features from several sources outweigh that.
% - The held-out outer folds answer this, not the inner search.
% - Chapter 6 computes the importance ranking on noisy labels, so it tests this worry directly.
% Paragraph 5: the budget and the information bottleneck
% - The information bottleneck finds a short code for the input that keeps as much information as possible about the target \citep{Tishby2000}.
% - \citet{Tishby2015} applies this to the layers of deep networks. It does not show that compression helps.
% - The budget is an analogy only. It keeps a subset of the original features; the bottleneck learns a compressed mapping of the input.
% - The evidence is the comparison at matched numbers of features, not the analogy.
% Paragraph 6: what a larger search buys (only if the search-size runs are made)
% - Question: does the hybrid's lead over the comparators grow with the number of configurations searched?
% - If the lead grows, the method's value depends on how much search it is given.
% - The comparators get no search. Read only the lead after leave-one-fold-out reselection, which charges the hybrid for taking the best of more draws.
% Paragraph 7: whether the internal evaluator must match the final model
% - The settings: matched, a near family, a far family, a wider network for the DNN, and a fixed set of three tree evaluators.
% - "Matched" means the same model family at the search's lighter settings, not the downstream model's exact settings. Say so.
% - If the settings score alike, the model scoring each representation during the search need not be the model finally trained. That is a finding, not a failed test.
% - One possible reason: if the features that carry signal are the same for every model family, any reasonable evaluator finds them.
% - This sits against the premise that different prediction tasks prefer different settings, and against what different downstream models pick. Say how they fit together.
% - Consequence: one cheap evaluator may be enough, so the search need not be re-run per model for this reason.
% Paragraph 8: where the hybrid does not help
% - On some datasets the best single representation scores as well as the hybrid, or better.
% - One explanation to test: one representation already holds most of the signal, so features from the others add little. Read it from that representation's share of the hybrid's features; do not state it as found.
% - A tie there is consistent with the search not inflating scores. A search that still reported a large gain there would be the warning sign of overfitting.
% - Practical signal: score each representation alone before searching. If one stands well clear of the others, a small gain is likely.
% - \citet{Laufktter2019}: ECFP4 joined with a fingerprint built from 561 PubChem screening assays, random forest, 24 held-out assays, sixfold cross-validation. The combination had higher ROC-AUC than both single fingerprints on 18 of the 24 assays. On the rest the paper calls it "comparable" to the better of the two.
% - Its MCC was higher than both on all but 5 of the 24 assays. So it reports ties with the best single fingerprint on a minority of assays; it does not say when to expect them.
% - \citet{Turkina2025}: six non-hashed fingerprints combined by low-, mid- and high-level fusion, six public datasets, three regression and three classification. The abstract says only that mid-level fusion "consistently improves predictive accuracy".
% - TODO(lit): neither source above is a study of when combining gives no gain. A source that reports it directly is still needed.
% ---- END NOTES ----

% ===== FROM: Limitations notes =====
% ---- NOTES: Limitations ----
% Takeaway: The search is costly, must be repeated when data, representations or models change, and its gain is unknown until it has run.
% Takeaway: On small datasets, frozen inputs and noisy endpoints the hybrid has a ceiling, and simpler selection may be enough.
% Paragraph 1: small datasets, frozen embeddings, sequence encodings, integration burden, noise ceiling (her paragraph below)
% - Change: "KIRBy's limitations" should read "the hybrid method's limitations". "Best single source" should read "best single representation".
% - The search scores each configuration on a held-out part of each training fold: one fifth of the training molecules, split by the same groups as the folds.
% - With 5 outer folds, a 200-molecule dataset leaves 160 training molecules and a held-out part of about 32 molecules. About 20 held-out molecules corresponds to about 125 molecules in the whole dataset. Change: match the two numbers in the prose.
% - TODO(lit): source for SchNet outperforming frozen embeddings on 3D quantum-mechanical properties. The sentence has no citation. The SchNet paper (Sch\"utt et al., NeurIPS 2017) and Faber et al. 2017 do not compare SchNet with frozen or pretrained embeddings.
% - Change: "20 to 50 foundation models" is true of the software as a whole, not of one search. The dataset lists draw on 33 pretrained models in total, across all domains.
% - One search sees far fewer. Most lists hold 10 base representations. On the small-molecule datasets 3 to 7 of the 10 are pretrained models. On the antibody datasets 6 to 10 are, but several are layers or regions of the same model, so 3 to 8 distinct models.
% - Change: reword the integration sentence to match: the burden is building and maintaining 33 pretrained models in one codebase; each search then uses a handful.
% - Change: one-hot and $k$-mer features are themselves base representations (one-hot on several antibody datasets, $k$-mer counts from length 3 to 6 on the DNA datasets). The sentence should not read as though they are left out. It needs a citation or a result for "disperses the positional and compositional information", or should be cut.
% - Change: \citet{Kramer2012} gives a standard deviation of 0.54 pK$_i$ units between independent measurements, and a maximum possible squared Pearson correlation of 0.81 on large datasets. Say "standard deviation of 0.54 pK$_i$ units" and "on large datasets".
% - Change: "Kramer et al.\ \citep{Kramer2012}" prints the names twice. Use \citet{Kramer2012}.
% - Change: "The budget constraint plausibly guards against overshoot": no result in the thesis measures this. Keep it as a guess in words, or cut it.
% - The budget is not capped at the number of training molecules. Budgets above that number are measured (selection bias against budget).
% - Result needed: whether the hybrid fails to beat the best single representation below about 200 molecules. Keep the claim only if the full runs show it.
% Paragraph 2: compute and run time
% - The search is a full grid. Every configuration is scored on every fold unless pruning is on. Under nested cross-validation the whole grid runs again inside every outer fold.
% - Importance scores are computed once per importance method per fold, on training data only. All allocation methods and budgets in that fold reuse them.
% - One exception: the global pool runs one extra importance fit per configuration.
% - Importance methods differ in cost. SHAP values and DeepLIFT attributions cost far more than tree gain importance.
% - Cost grows with the number of molecules (every fit), with the total number of features across representations (importance over every feature) and with the number of representations.
% - Caching repeated work (fold splits, importance scores, each representation's own score) cuts the cost of a re-run. It does not cut the cost of the first run.
% - Consequence: the search costs orders of magnitude more than fitting one model on one representation. It needs a cluster, which limits who can use it.
% - Figure: run time (hours, log scale) against number of molecules, full search. One point per dataset, coloured by domain. Horizontal reference line: time to fit the best single representation once.
% - Result needed: hours per dataset for the full search, and which importance method and allocation method take most of the time.
% Paragraph 3: the search must be repeated when inputs change
% - A search result is fitted to one dataset, one set of base representations, one downstream model and one set of folds.
% - New molecules or different folds change the importance scores and each representation's score. The budget split and the selected features can then change.
% - A new representation: each representation's share of the budget is set relative to the others. The whole search reruns, not only the new representation's part.
% - A new downstream model: the best configuration is chosen per model, so a new model needs its own search.
% - This is the same kind of cost as re-tuning hyperparameters, over a larger space.
% - Consequence: datasets grow (new assay rounds, new chemical series), and each update means paying for the search again.
% - Partial workaround: reusing a blueprint on a related dataset. Its limits belong here.
% - Result needed: how far a blueprint carries to new folds or a related dataset (overlap of selected features, change in score in the dataset's own metric).
% Paragraph 4: high upfront cost, gain unknown until the search has run
% - No cheap test before the search predicts whether the hybrid will beat the best single representation.
% - Some datasets are not helped, typically when one representation dominates. A standout representation in the pool is a hint, not a guarantee.
% - Consequence: on those datasets the full cost buys the answer that the single representation is enough.
% - Figure: per dataset, hybrid score minus best single-representation score (dataset's own metric) against search cost in hours. Per-fold points, no error bars.
% Paragraph 5: reproducibility
% - A hybrid depends on many settings: the representations and their versions, the fold seed, the importance method, the allocation method, the budget, the internal evaluator, and library versions.
% - Small changes to the training molecules can change which features are selected. The stability index measures this.
% - Some importance methods are stochastic (DeepLIFT attributions, subsampled tree ensembles). Rebuilding the same hybrid needs the same seed and library versions.
% - A blueprint stores the settings, not the selected base-representation features. A rebuild on other molecules can pick other features.
% - Pretrained embeddings add a dependence on the checkpoint version. A re-released checkpoint changes the representation.
% - Consequence: reproducing a number needs the full set of settings, not "a hybrid of ECFP and descriptors" (Appendix~\ref{app:availability}; Additional files, Section~\ref{app:searchspace}).
% Paragraph 6: simpler selection may be enough
% - Simpler options: the best single representation chosen by cross-validation; concatenation with a filter or an $L_1$ penalty; a tree ensemble's built-in selection. Each costs about one fit or one short tuning run.
% - Any of them tuned by cross-validation also needs nested cross-validation for an honest estimate \citep{Varma2006}. The cost gap is smaller than it first looks.
% - Rule to state: with limited compute or a routine task, start with the best single representation. Run the search when a gain beyond the fold-to-fold spread would matter.
% Paragraph 7: scope of the claims
% - The grid is finite. A better hybrid may exist at other budgets or with other importance methods; the search finds the best configuration within the grid only.
% - Close: these limits make the method a tool for choosing features when a gain is worth the compute, not a default for every dataset.
% ---- END NOTES ----

% ===== FROM: Conclusions notes =====
% ---- NOTES: Conclusions ----
% Takeaway: One paragraph per question this chapter answers, then a short cost line, then the bridge to Chapter 6.
% - Quote each question's wording from Section~\ref{sec:intro:questions} at the start of its paragraph. State the answer in one or two sentences.
% Paragraph 1: how do you select the best set of features for a task? (does the hybrid beat its inputs)
% - Answer against two comparators: the best single representation at its full width, and an ensemble of one model per representation.
% - Framing: where the hybrid only matches the best single representation, the search still found that representation without it being named in advance.
% - The premise tested here, argued in Chapter 4: different representations carry different information.
% - Result needed: hybrid against each comparator, per dataset, in the dataset's own metric, with per-fold spread.
% - Result needed: the datasets where the hybrid did not help, and which single representation dominated on each.
% Paragraph 2: selection, not only reduction
% - At the same number of features, the hybrid is compared against PCA, mRMR, recursive feature elimination and random features.
% - If the hybrid is ahead at matched width, the gain comes from which features are chosen, not from having fewer features.
% - The evidence is the comparison at matched numbers of features, not the information-bottleneck analogy.
% - Result needed: hybrid minus each comparator at matched width, per dataset.
% - Result needed: best budget per downstream model (features), and whether tree ensembles and neural networks prefer different widths.
% - Result needed (only if the search-size comparison is kept): hybrid minus each comparator at a small and a large search, after leave-one-fold-out reselection.
% Paragraph 3: one search setting is not enough
% - Every searched setting can pick a different value on different datasets: allocation method, number of representations, importance method, internal evaluator, downstream model, budget.
% - Consequence: a fixed recipe (for example, top two representations, one importance method, one budget) would lose score on some datasets. This is why the search exists.
% - Framing only, not proof: no single optimiser is best when averaged over all problems \citep{Wolpert1997}.
% - Result needed: per setting, the number of datasets on which each value was chosen.
% Paragraph 4: is the search overfitting?
% - The whole search runs inside each outer fold, and importance is computed on training molecules only (nested cross-validation). One sentence.
% - Result needed: selection bias per budget (features), and the stability index of the selected features per dataset.
% Paragraph 5: what blueprints tell a chemist
% - Descriptor features and unhashed substructure keys have chemical names. Features from learned embeddings do not.
% - Result needed: one or two datasets where the named features that were selected match known chemistry of the endpoint.
% Paragraph 6: beyond small molecules, transfer, and data held at separate sites
% - One sentence on what a dataset needs for the hybrid to help: enough samples, and base representations whose features carry importance.
% - Label noise sets a separate ceiling on any gain (the public $K_i$ uncertainty estimate).
% - Result needed: antibodies, reactions, nucleotides and gene expression: hybrid against the best single representation on each.
% - Result needed: transfer across datasets, the antibody mutational scans, and data held at separate sites. One sentence each on the outcome.
% Paragraph 7: are hybrids more robust to label noise than their base representations? (only if this comparison stays in this chapter)
% - Result needed: \aucnorm, hybrid at each budget against each single representation, per dataset and noise type, with per-fold spread.
% Paragraph 8: does importance weighting help? (the parts this chapter answers)
% - Here weighting means using importance scores to scale features or molecules during training, not only to rank and select features.
% - Five training conditions on the hybrid: unweighted, features scaled by importance, molecules weighted by agreement with their neighbours, both, and only the features above the median importance.
% - Caveat: the configuration comes from a search over all molecules; only the feature selection is redone per fold.
% - Weighting in distances and under learning with noisy labels is answered in Chapter 6.
% - Result needed: each weighted condition minus unweighted, per dataset and model, in the dataset's own metric.
% Paragraph 9: cost, stated once
% - One or two sentences: the search is expensive, and it must be run again for new folds, a new representation or a new downstream model.
% - Simpler feature selection may be enough on some tasks.
% Paragraph 10: bridge to Chapter 6
% - Chapter 4 argued that a fixed fingerprint with Tanimoto similarity spends most of its features on information unrelated to the task.
% - This chapter used the hybrid only as input to a model.
% - Chapter 6 uses it for similarity, in two ways, and asks whether either improves noise detection and learning with noisy labels, which Chapter 4 showed failing under standard similarity.
% - Why it matters: if a feature set chosen for the task gives better neighbours, every method that relies on neighbours could use it. Examples: noise detection, applicability domain, nearest-neighbour prediction.
% - Figure (optional): three boxes in a row. Chapter 4 (standard representations and similarity fail) to Chapter 5 (a feature set chosen for the task) to Chapter 6 (that feature set used for similarity and label noise).
% ---- END NOTES ----
